\documentclass[12pt]{article}

%% ── Packages ────────────────────────────────────────────────────────────────
\usepackage{amsmath, amssymb, amsthm}
\usepackage{geometry}
\usepackage{hyperref}
\usepackage{bm}
\usepackage{xcolor}
\usepackage{tcolorbox}
\usepackage{booktabs}

\geometry{margin=1.1in}
\hypersetup{colorlinks=true, linkcolor=blue, citecolor=blue}

%% ── Custom environments ─────────────────────────────────────────────────────
\tcbuselibrary{breakable, skins}

\newtcolorbox{derivbox}[1][]{
  breakable, enhanced,
  colback=blue!4, colframe=blue!40!black,
  fonttitle=\bfseries,
  title={#1}
}

\newtcolorbox{notebox}[1][Note]{
  breakable, enhanced,
  colback=orange!6, colframe=orange!60!black,
  fonttitle=\bfseries,
  title={#1}
}

\newtcolorbox{keybox}[1][Key Result]{
  breakable, enhanced,
  colback=green!5, colframe=green!50!black,
  fonttitle=\bfseries,
  title={#1}
}

%% ── Macros ──────────────────────────────────────────────────────────────────
\newcommand{\dd}{\mathrm{d}}
\newcommand{\ii}{\mathrm{i}}
\newcommand{\half}{\tfrac{1}{2}}
\newcommand{\bu}{\bm{u}}
\newcommand{\br}{\bm{r}}
\renewcommand{\vec}[1]{\bm{#1}}

%% ── Theorem styles ──────────────────────────────────────────────────────────
\theoremstyle{definition}
\newtheorem{definition}{Definition}[section]
\newtheorem{result}{Result}[section]

%% ────────────────────────────────────────────────────────────────────────────
\title{\textbf{Oscillations of a Viscous Liquid Drop}\\[6pt]
       \large A Self-Contained Pedagogical Expansion of Reid (1960)}
\author{Expanded from W.\,H.\ Reid, \emph{Quart.\ Appl.\ Math.}\ \textbf{18}(1), 86--89 (1960)}
\date{}

\begin{document}
\maketitle
\tableofcontents
\bigskip

%% ════════════════════════════════════════════════════════════════════════════
\section{Introduction and Physical Motivation}
%% ════════════════════════════════════════════════════════════════════════════

\subsection{What problem are we solving?}

Consider a small liquid droplet (say, a water drop suspended in air) that has been
slightly deformed from its equilibrium spherical shape. Surface tension acts as a
restoring force, and the drop will oscillate. Viscosity damps those oscillations. The
question is: \emph{at what frequency does it oscillate, and how fast does it decay?}

This couples three distinct pieces of physics:
\begin{enumerate}
  \item The inviscid capillary (surface-tension-driven) oscillation modes, characterized
        by the frequency $\sigma_{l;0}$ for a deformation of spherical-harmonic order $l$.
  \item Viscous dissipation inside the bulk of the drop.
  \item The coupling between surface deformation, internal flow, and the pressure
        field, all of which must be consistent at the boundary.
\end{enumerate}

Reid's paper resolves all three simultaneously for \emph{arbitrary} viscosity, by reducing
the problem to a single transcendental characteristic equation. Its central insight is a rescaling. Surface tension enters only through the inviscid
frequency $\sigma_{l;0}$. After non-dimensionalisation, the characteristic equation is
identical to Chandrasekhar's equation for a self-gravitating viscous liquid globe.

\subsection{Historical context and relation to Mola\v{c}ek \& Bush (2012)}

Three analytical milestones precede Reid's result:
\begin{itemize}
  \item \textbf{Rayleigh (1879)} / \textbf{Lamb (1932)}: inviscid frequencies, Eq.~\eqref{eq:inviscid_freq}.
  \item \textbf{Lamb (1881)}: small-viscosity correction,
        $\sigma_{l;\nu} = (l-1)(2l+1)\nu/R^2 \pm \ii\,\sigma_{l;0}$.
  \item \textbf{Chandrasekhar (1959)}: full arbitrary-viscosity solution for a
        \emph{self-gravitating} globe.
\end{itemize}
Reid (1960) shows the surface-tension problem is identical to Chandrasekhar's,
thereby completing the picture for surface-tension-driven drops.

The result is directly used in Mola\v{c}ek \& Bush (2012), where the coefficients
$A_m$ and $D_m$ governing kinetic energy and viscous dissipation for each surface mode
are extracted from Reid's characteristic equation (Eq.~\eqref{eq:char_eq} below).

%% ════════════════════════════════════════════════════════════════════════════
\section{Mathematical Preliminaries}
%% ════════════════════════════════════════════════════════════════════════════

\subsection{Spherical harmonics}

A \emph{spherical harmonic} of degree $l$ and order $m$ is the function
\begin{equation}
  Y_l^m(\theta,\varphi) = P_l^m(\cos\theta)\,e^{-\ii m\varphi},
  \label{eq:sph_harm}
\end{equation}
where $P_l^m$ is the associated Legendre polynomial, $\theta$ is the polar angle from
the $z$-axis, and $\varphi$ is the azimuthal angle. (We use the unnormalized convention;
the characteristic equation is independent of the normalization choice.)

The key property we use is that spherical harmonics are \emph{eigenfunctions of the
Laplace--Beltrami operator} (the angular part of the Laplacian):
\begin{equation}
  \nabla^2_{\text{angular}}\,Y_l^m
  = -\frac{l(l+1)}{r^2}\,Y_l^m.
  \label{eq:sph_harm_eigen}
\end{equation}
This means that for any function of the form $f(r)\,Y_l^m(\theta,\varphi)$, the full
scalar Laplacian is
\begin{equation}
  \nabla^2\!\left[f(r)\,Y_l^m\right]
  = \left[\,f'' + \frac{2}{r}f' - \frac{l(l+1)}{r^2}\,f\,\right] Y_l^m,
  \label{eq:scalar_lap}
\end{equation}
where primes denote $\dd/\dd r$.

For axisymmetric problems we take $m=0$, so $Y_l^0 \propto P_l(\cos\theta)$ is just the
Legendre polynomial. Reid's analysis holds for all $m$; the characteristic equation
turns out to be independent of $m$.

\subsection{Spherical Bessel functions}
\label{sec:bessel}

The ordinary Bessel equation of order $\nu$ is
\begin{equation}
  w'' + \frac{1}{z}w' + \left(1 - \frac{\nu^2}{z^2}\right)w = 0,
\end{equation}
with solutions $J_\nu(z)$ (Bessel functions of the first kind) and $Y_\nu(z)$
(second kind, singular at $z=0$).

The \emph{spherical Bessel equation} of order $l$ is
\begin{equation}
  v'' + \frac{2}{z}v' + \left(1 - \frac{l(l+1)}{z^2}\right)v = 0,
  \label{eq:sph_bessel}
\end{equation}
and arises when separating the Helmholtz equation $(\nabla^2 + k^2)F = 0$ in spherical
coordinates. Its regular solution is the \emph{spherical Bessel function of the first kind},
\begin{equation}
  j_l(z) = \sqrt{\frac{\pi}{2z}}\,J_{l+1/2}(z).
\end{equation}

\begin{notebox}
Any function $U(x) = x\,j_l(qx)$ satisfies the ODE
\begin{equation}
  U'' - \frac{l(l+1)}{x^2}\,U + q^2 U = 0
  \label{eq:helm_U}
\end{equation}
(the prime here is $\dd/\dd x$). This follows immediately by the substitution $U = xv$
with $x = qr/R$, which reduces \eqref{eq:helm_U} to the spherical Bessel equation
\eqref{eq:sph_bessel} for $v$.

\begin{proof}
Let $U = xv(x)$. Then $U' = v + xv'$ and $U'' = 2v' + xv''$. Substituting:
\begin{align*}
  2v' + xv'' - \frac{l(l+1)}{x}\,v + q^2 xv
  &= x\!\left[v'' + \frac{2}{x}v' + q^2 v - \frac{l(l+1)}{x^2}v\right] = 0,
\end{align*}
which is precisely \eqref{eq:sph_bessel} with $z = qx$.
\end{proof}
\end{notebox}

We will also need the derivative identity. Using the standard recurrence
$j_l'(z) = \frac{l}{z}j_l(z) - j_{l+1}(z)$, we obtain
\begin{equation}
  \frac{q\,j_l'(q)}{j_l(q)} = l - q\,\frac{j_{l+1}(q)}{j_l(q)}.
  \label{eq:bessel_ratio}
\end{equation}
Reid defines the ratio
\begin{equation}
  Q_{l+1/2}(q) \equiv \frac{J_{l+3/2}(q)}{J_{l+1/2}(q)} = \frac{j_{l+1}(q)}{j_l(q)},
  \label{eq:Q_def}
\end{equation}
so \eqref{eq:bessel_ratio} becomes $q j_l'(q)/j_l(q) = l - q Q_{l+1/2}(q)$.

\subsection{The poloidal decomposition}
\label{sec:poloidal}

Any divergence-free vector field $\bu$ can be uniquely decomposed into a
\emph{toroidal} part (involving only $\nabla\times(\Psi\hat{r})$, with no radial
component) and a \emph{poloidal} part (involving $\nabla\times\nabla\times(\Phi\hat{r})$,
with non-zero radial component). For axisymmetric flow with no swirl, the toroidal
part vanishes, and $\bu$ is purely poloidal.

For a poloidal field with angular dependence $Y_l^m$, the entire velocity field is
determined by its radial component $u_r$. Incompressibility then gives $u_\theta$,
and there are no other free components. This is why it suffices to write down a
single scalar ODE for $U(x)$, where $u_r \propto U(x)/x^2$.


%% ════════════════════════════════════════════════════════════════════════════
\section{Problem Setup}
%% ════════════════════════════════════════════════════════════════════════════

\subsection{Equilibrium state}

A liquid drop of density $\rho$, viscosity $\mu = \rho\nu$, and surface tension $T_1$
is in equilibrium as a sphere of radius $R$. With the external pressure set to zero,
the \emph{Young--Laplace equation} gives the internal equilibrium pressure:
\begin{equation}
  p_0 = \frac{2T_1}{R}.
  \label{eq:YL_equilibrium}
\end{equation}

\begin{notebox}[Young--Laplace Equation]
For an interface with principal radii of curvature $R_1$ and $R_2$, the pressure
jump across the interface is $\Delta p = T_1(1/R_1 + 1/R_2)$. For a sphere,
$R_1 = R_2 = R$, giving $\Delta p = 2T_1/R$.
\end{notebox}

\subsection{Perturbed surface}

We consider a small deformation of the spherical surface of the form
\begin{equation}
  r = R\!\left[1 + \epsilon\,Y_l^m(\theta,\varphi)\right],
  \qquad \epsilon \ll 1.
  \label{eq:surface_deform}
\end{equation}
We look for solutions with time dependence
\begin{equation}
  \epsilon(t) = \epsilon_0\,e^{-\sigma t},
  \label{eq:time_dep}
\end{equation}
where $\sigma$ is a complex number to be determined. The convention is chosen so that
$\mathrm{Re}(\sigma) > 0$ corresponds to decay.

\begin{notebox}
Why $e^{-\sigma t}$ rather than $e^{+\ii\omega t}$? Reid follows the convention of
Chandrasekhar, where $\sigma$ is complex and damped oscillations correspond to
$\sigma = \gamma + \ii\omega_d$ with $\gamma > 0$ (decay) and $\omega_d > 0$
(oscillation frequency). In the inviscid limit, $\gamma\to 0$ and $\sigma\to \pm\ii\omega_d$.
\end{notebox}

\subsection{Inviscid frequencies}

For completeness, we record Lamb's result for the oscillation frequencies in the
\emph{absence} of viscosity. The frequency of a spherical harmonic mode of order $l$ is
\begin{equation}
  \sigma_{l;0}^2 = l(l-1)(l+2)\,\frac{T_1}{\rho R^3}.
  \label{eq:inviscid_freq}
\end{equation}
These are real (undamped oscillations) and independent of the azimuthal order $m$.
The $l=2$ mode, for example, corresponds to the oblate--prolate oscillation familiar
from a wobbling drop of water.

\begin{notebox}[Physical check]
Equation~\eqref{eq:inviscid_freq} is dimensionally consistent:
$[T_1/(\rho R^3)] = [\text{N/m}]/[\text{kg/m}^3 \cdot \text{m}^3]
 = [\text{kg/s}^2]/[\text{kg}] = [\text{s}^{-2}]$. The $l=1$ mode gives
$\sigma_{1;0}^2 = 0$, corresponding to a rigid translation (no restoring force).
\end{notebox}

%% ════════════════════════════════════════════════════════════════════════════
\section{Linearized Governing Equations}
%% ════════════════════════════════════════════════════════════════════════════

The departures from equilibrium are governed by the linearized incompressible
Navier--Stokes equations,
\begin{equation}
  \frac{\partial\bu}{\partial t} = -\nabla\frac{\delta p}{\rho}
                                    - \nu\,\nabla\times\nabla\times\bu,
  \qquad
  \nabla\cdot\bu = 0.
  \label{eq:NS_lin}
\end{equation}

\begin{notebox}[Vector identity]
For incompressible flow, $\nabla\cdot\bu = 0$ implies
$\nabla^2\bu = -\nabla\times\nabla\times\bu$, so
\eqref{eq:NS_lin} is equivalent to
$\partial_t\bu = -\nabla(\delta p/\rho) + \nu\nabla^2\bu$.
Both forms appear in the literature.
\end{notebox}

With the time dependence \eqref{eq:time_dep}, all fields vary as $e^{-\sigma t}$,
so $\partial_t \to -\sigma$. The momentum equation becomes
\begin{equation}
  -\sigma\bu = -\nabla\frac{\delta p}{\rho} + \nu\nabla^2\bu.
  \label{eq:NS_sigma}
\end{equation}

%% ════════════════════════════════════════════════════════════════════════════
\section{The Pressure Field}
%% ════════════════════════════════════════════════════════════════════════════

\subsection{Pressure satisfies Laplace's equation}

\begin{derivbox}[Derivation: $\nabla^2(\delta p/\rho) = 0$]
Take the divergence of \eqref{eq:NS_sigma}:
\[
  -\sigma\underbrace{(\nabla\cdot\bu)}_{=\,0} 
  = -\nabla^2\!\left(\frac{\delta p}{\rho}\right) 
  + \nu\underbrace{\nabla\cdot(\nabla^2\bu)}_{\nabla^2(\nabla\cdot\bu)\,=\,0},
\]
where we used $\nabla\cdot\nabla^2\bu = \nabla^2(\nabla\cdot\bu) = 0$.
Therefore:
\begin{equation}
  \nabla^2\!\left(\frac{\delta p}{\rho}\right) = 0.
  \label{eq:pressure_Laplace}
\end{equation}
\end{derivbox}

The pressure perturbation satisfies Laplace's equation. We need the solution that:
(a) has angular dependence $Y_l^m$, and (b) is \emph{regular at the origin} $r=0$.

From \eqref{eq:scalar_lap}, the radial ODE for $f(r)$ in $\nabla^2(f\,Y_l^m) = 0$~is
\begin{equation}
  f'' + \frac{2}{r}f' - \frac{l(l+1)}{r^2}f = 0.
\end{equation}
The general solution is $f(r) = A r^l + B r^{-(l+1)}$. Regularity at $r=0$ requires
$B = 0$. Writing $x = r/R$, we therefore have
\begin{equation}
  \frac{\delta p}{\rho} = \epsilon\,P_0\,x^l\,Y_l^m,
  \label{eq:pressure_soln}
\end{equation}
where $P_0$ is a constant of integration (with dimensions of velocity squared)
to be determined from the boundary conditions.

\subsection{Scalar potential for the pressure gradient}

Since $\nabla(\delta p/\rho)$ is a purely poloidal, divergence-free vector field
with angular structure $Y_l^m$, it can be represented through a single scalar function
$\Pi(x)$ via its radial component:
\begin{equation}
  \left(\nabla\frac{\delta p}{\rho}\right)_r
  = \epsilon_0\,\sigma^2 R\,\frac{\Pi(x)}{x^2}\,Y_l^m\,e^{-\sigma t}.
  \label{eq:Pi_def}
\end{equation}
Computing the radial component of $\nabla(\epsilon P_0 x^l Y_l^m)$:
\[
  \left(\nabla\frac{\delta p}{\rho}\right)_r
  = \frac{\partial}{\partial r}\!\left(\epsilon P_0 x^l Y_l^m\right)
  = \epsilon P_0\,\frac{l\,x^{l-1}}{R}\,Y_l^m
  = \epsilon_0 e^{-\sigma t} P_0\,\frac{l}{R}\,x^{l-1}\,Y_l^m.
\]
Comparing with \eqref{eq:Pi_def}, we need
$\epsilon_0\sigma^2 R\,\Pi(x)/x^2 = \epsilon_0 P_0 l x^{l-1}/R$,
which gives $\Pi(x) = \Pi_0\,x^{l+1}$ with
\begin{equation}
  \Pi(x) = \Pi_0\,x^{l+1}, \qquad
  \Pi_0 = \frac{l}{\sigma^2 R^2}\,P_0.
  \label{eq:Pi_form}
\end{equation}

%% ════════════════════════════════════════════════════════════════════════════
\section{The Velocity Field and its Governing ODE}
\label{sec:velocity_ode}
%% ════════════════════════════════════════════════════════════════════════════

\subsection{Poloidal representation}

Since the pressure gradient is purely poloidal, the momentum equation
\eqref{eq:NS_sigma} forces $\bu$ to also be purely poloidal (the toroidal part would
have no driver). For a poloidal field with angular dependence $Y_l^m$, we write the
radial component as
\begin{equation}
  u_r = \epsilon_0\,\sigma R\,\frac{U(x)}{x^2}\,Y_l^m\,e^{-\sigma t},
  \label{eq:ur_def}
\end{equation}
where $U(x)$ is a dimensionless scalar function.

\begin{notebox}[Why $U(x)/x^2$?]
The scaling $1/x^2$ is conventional and simplifies the algebra.
It arises from the stream-function representation of poloidal flow, and ensures
the ODE for $U$ takes the standard form \eqref{eq:ODE_U} below.
\end{notebox}

\subsection{Derivation of the ODE for $U(x)$}

\begin{derivbox}[Full derivation]
We substitute the ansatz $u_r = \sigma R\,G(x)\,Y_l^m$ (where $G = U/x^2$, factoring out
the $\epsilon_0 e^{-\sigma t}$ terms throughout) into the $r$-component of
\eqref{eq:NS_sigma}. The $r$-component of the vector Laplacian for an axisymmetric
poloidal field is:
\begin{equation}
  \left[\nabla^2\bu\right]_r = \nabla^2 u_r + \frac{2}{r^2}u_r + \frac{2}{r}\frac{\partial u_r}{\partial r}.
  \label{eq:vec_lap_r}
\end{equation}

\textit{Derivation of \eqref{eq:vec_lap_r}.}
The vector Laplacian r-component in spherical coordinates is:
\[
  \left[\nabla^2\bu\right]_r = \nabla^2 u_r - \frac{2u_r}{r^2}
    - \frac{2}{r^2\sin\theta}\frac{\partial(\sin\theta\,u_\theta)}{\partial\theta}.
\]
From incompressibility, $(1/r\sin\theta)\,\partial(\sin\theta\,u_\theta)/\partial\theta
= -(1/r^2)\,\partial(r^2 u_r)/\partial r$, so:
\[
  \frac{1}{\sin\theta}\frac{\partial(\sin\theta\,u_\theta)}{\partial\theta}
  = -\frac{r}{r^2}\,\partial_r(r^2 u_r)
  = -(2u_r + r u_r').
\]
Substituting:
\[
  \left[\nabla^2\bu\right]_r 
  = \nabla^2 u_r - \frac{2u_r}{r^2} 
    + \frac{2}{r^2}(2u_r + r u_r')
  = \nabla^2 u_r + \frac{2u_r}{r^2} + \frac{2}{r}u_r'.
\]

Now, for $u_r = \sigma R\,G(x)\,Y_l^m$ with $G = G(x)$ and $x=r/R$:
\[
  \nabla^2 u_r = \sigma R\,Y_l^m\left[G'' + \frac{2}{x}G' - \frac{l(l+1)}{x^2}G\right]\frac{1}{R^2}
\quad\text{(from \eqref{eq:scalar_lap})},
\]
\[
  \frac{2}{r}u_r' = \frac{2}{Rx}\cdot \sigma\,G'\,Y_l^m
  = \frac{2\sigma}{R x}\,G'\,Y_l^m,
\]
\[
  \frac{2u_r}{r^2} = \frac{2\sigma R G}{R^2 x^2}\,Y_l^m = \frac{2\sigma G}{R x^2}\,Y_l^m.
\]
Combining, the $r$-component of the momentum equation $-\sigma\bu = -\nabla(\delta p/\rho) + \nu\nabla^2\bu$ becomes (after canceling $Y_l^m e^{-\sigma t}$ throughout and using $x = r/R$):
\begin{equation}
  -\sigma^2 R\,G = -\frac{P_0 l\,x^{l-1}}{R} + \frac{\nu\sigma}{R}\left[\,G'' + \frac{4}{x}G' + \frac{2 - l(l+1)}{x^2}G\,\right].
  \label{eq:r_comp_mom}
\end{equation}

Multiplying through by $R/(\nu\sigma)$ and using $q^2 = \sigma R^2/\nu$:
\[
  -q^2 G = -\frac{P_0 l x^{l-1}}{\sigma\nu} + G'' + \frac{4}{x}G' + \frac{2-l(l+1)}{x^2}G.
\]

Now substitute $G = U/x^2$. One computes:
\[
  G'' + \frac{4}{x}G' + \frac{2-l(l+1)}{x^2}G 
  = \frac{U'' - l(l+1)U/x^2}{x^2}.
\]
(This identity is most easily verified by direct substitution, expanding each term.)

Multiplying through by $x^2$:
\[
  U'' - \frac{l(l+1)}{x^2}U + q^2 U = \frac{P_0 l x^{l+1}}{\sigma\nu}.
\]
Using \eqref{eq:Pi_form}, $P_0 = \sigma^2 R^2 \Pi_0/l$, so
$P_0 l/(\sigma\nu) = \sigma R^2 \Pi_0/\nu = q^2\Pi_0$. Therefore:
\begin{equation}
  \boxed{
    \left[\frac{\dd^2}{\dd x^2} - \frac{l(l+1)}{x^2} + q^2\right]U(x) = q^2\Pi_0\,x^{l+1},
  }
  \label{eq:ODE_U}
\end{equation}
where
\begin{equation}
  q^2 = \frac{\sigma R^2}{\nu}.
  \label{eq:q_def}
\end{equation}
This is Reid's Eq.~(9). \hfill$\square$
\end{derivbox}

\subsection{Solving the ODE}

Equation~\eqref{eq:ODE_U} is a second-order linear inhomogeneous ODE. We find
the general solution by superposing a particular solution and the general
homogeneous solution.

\paragraph{Particular solution.}
Try $U_p = \Pi_0\,x^{l+1}$:
\[
  U_p'' = l(l+1)\Pi_0\,x^{l-1}, 
  \quad
  \frac{l(l+1)}{x^2}U_p = l(l+1)\Pi_0\,x^{l-1},
  \quad
  q^2 U_p = q^2\Pi_0\,x^{l+1}.
\]
The first two terms cancel exactly, leaving $q^2\Pi_0\,x^{l+1}$ on the left, matching
the right-hand side. So
\begin{equation}
  U_p(x) = \Pi_0\,x^{l+1}.
\end{equation}

\paragraph{Homogeneous solution.}
The homogeneous equation $U'' - l(l+1)U/x^2 + q^2 U = 0$ has solutions $x\,j_l(qx)$
and $x\,n_l(qx)$, where $j_l$ and $n_l$ are spherical Bessel functions of the first
and second kind (see Section~\ref{sec:bessel}). Since $n_l(z) \sim -(2l-1)!!\,z^{-(l+1)}$ as
$z\to 0$ (diverges at the origin), regularity requires we discard $n_l$. The regular homogeneous
solution is thus $C\,x\,j_l(qx)$.

\paragraph{General solution.}
\begin{equation}
  U(x) = C\,x\,j_l(qx) + \Pi_0\,x^{l+1},
  \label{eq:U_general}
\end{equation}
where $C$ and $\Pi_0$ are determined by the boundary conditions.

%% ════════════════════════════════════════════════════════════════════════════
\section{Boundary Conditions}
%% ════════════════════════════════════════════════════════════════════════════

Three boundary conditions must be imposed at the drop surface $x = 1$.

\subsection{BC1: Kinematic condition}
\label{sec:bc1}

The radial velocity of the fluid at the surface must equal the velocity at which
the surface itself moves.

\begin{derivbox}[Derivation of BC1]
The surface is at $r = R[1 + \epsilon\,Y_l^m]$ with $\epsilon = \epsilon_0 e^{-\sigma t}$.
The radial velocity of a surface element is:
\[
  \left.\frac{\partial r}{\partial t}\right|_{\text{surface}}
  = -\sigma R\,\epsilon_0 e^{-\sigma t}\,Y_l^m.
\]
The fluid radial velocity at $x = 1$ (from \eqref{eq:ur_def}) is:
\[
  u_r\big|_{x=1} = \epsilon_0\,\sigma R\,U(1)\,Y_l^m\,e^{-\sigma t}.
\]
Setting these equal:
\[
  \epsilon_0\,\sigma R\,U(1) = -\sigma R\,\epsilon_0.
\]
Canceling $\epsilon_0\sigma R$:
\begin{equation}
  U(1) = -1.
  \label{eq:BC1}
\end{equation}
The kinematic condition is imposed at $x=1$ (the unperturbed surface) because
evaluating at the actual surface $r=R[1+\epsilon Y_l^m]$ introduces only
$O(\epsilon^2)$ corrections.
\end{derivbox}

\begin{notebox}
The sign convention $U(1)=-1$ follows from the positive-$\sigma$ prefactor in
\eqref{eq:ur_def} with $\epsilon_0>0$ encoding the outward displacement.
This matches Reid (1960): his Eq.\,(8) uses the same positive-$\sigma R$ prefactor
and his Eq.\,(11) states $U(1)=-1$.
\end{notebox}

\subsection{BC2: Tangential stress condition}
\label{sec:bc2}

At a free surface with no viscous exterior fluid, the tangential viscous stress must
vanish.

The $r\theta$-component of the viscous stress tensor in spherical coordinates is:
\[
  \tau_{r\theta} = \mu\left[\,r\frac{\partial}{\partial r}\!\left(\frac{u_\theta}{r}\right) + \frac{1}{r}\frac{\partial u_r}{\partial\theta}\,\right].
\]
For a poloidal velocity field with harmonic structure $Y_l^m$, using the stream-function
representation and the incompressibility constraint to express $u_\theta$ in terms of $u_r$,
one can show that $\tau_{r\theta} = 0$ at
$x=1$ reduces to:
\begin{equation}
  \left[\,\frac{\dd^2}{\dd x^2} - \frac{2}{x}\frac{\dd}{\dd x} + \frac{l(l+1)}{x^2}\,\right]U = 0
  \qquad \text{at } x = 1.
  \label{eq:BC2}
\end{equation}

\begin{notebox}[Evaluating BC2 on the general solution]
For the particular solution $U_p = \Pi_0 x^{l+1}$:
\[
  \mathcal{L}_2[x^{l+1}] = l(l+1)x^{l-1} - 2(l+1)x^{l-1} + l(l+1)x^{l-1}
  = 2(l-1)(l+1)x^{l-1}.
\]
At $x=1$: $\mathcal{L}_2[U_p]|_{x=1} = 2(l-1)(l+1)\Pi_0$.

For the homogeneous part $U_h = Cxj_l(qx)$, using $U_h'' = -q^2 xj_l + l(l+1)j_l/x$
(from the Bessel equation) and $U_h' = j_l + qxj_l'$:
\[
  \mathcal{L}_2[U_h]\big|_{x=1}
  = C\!\left[-q^2 j_l(q) + 2\bigl(l(l+1)-1\bigr)j_l(q) - 2q j_l'(q)\right].
\]
Setting the sum to zero gives one equation relating $C$ and $\Pi_0$.
\end{notebox}

\subsection{BC3: Normal stress condition}

This boundary condition couples pressure, viscous stress, and surface curvature through surface tension.

\subsubsection{Mean curvature of the deformed surface}

\begin{derivbox}[Curvature of the perturbed sphere]
For the surface \eqref{eq:surface_deform}, the principal curvatures can be computed
by standard differential geometry. To first order in $\epsilon$:
\begin{equation}
  \frac{1}{R_1} + \frac{1}{R_2}
  = \frac{1}{R}\!\left[\,2 + (l-1)(l+2)\,\epsilon\,Y_l^m\,\right].
  \label{eq:curvature}
\end{equation}

\textit{Physical interpretation:} The $2/R$ term is the equilibrium curvature.
The perturbation contributes $(l-1)(l+2)\epsilon Y_l^m/R$. The factor $(l-1)(l+2)$
is the same one that appears in the inviscid frequency \eqref{eq:inviscid_freq}.
Both arise from the curvature response to a harmonic
deformation.
\end{derivbox}

\subsubsection{Normal stress balance}

The radial normal stress inside the drop is:
\begin{equation}
  -p_{rr} = p + \delta p - 2\mu\frac{\partial u_r}{\partial r}.
  \label{eq:normal_stress}
\end{equation}
The condition at the free surface is:
\begin{equation}
  -p_{rr} = T_1\!\left(\frac{1}{R_1} + \frac{1}{R_2}\right).
  \label{eq:normal_BC}
\end{equation}

\begin{derivbox}[Derivation of BC3]
At $O(\epsilon^0)$: the equilibrium Young--Laplace condition $p_0 = 2T_1/R$
is automatically satisfied (the drop is in equilibrium). At $O(\epsilon^1)$:
\[
  \delta p - 2\mu\frac{\partial u_r}{\partial r}\bigg|_{r=R}
  = T_1\,\frac{(l-1)(l+2)}{R}\,\epsilon\,Y_l^m.
\]
Using \eqref{eq:pressure_soln} with $x = 1$: $\delta p|_{x=1} = \rho\epsilon P_0 Y_l^m$
(since $\delta p/\rho = \epsilon P_0 x^l Y_l^m$ at $x=1$).

For the viscous term, with $u_r = \epsilon_0\sigma R\,G(x)Y_l^m e^{-\sigma t}$
and $G = U/x^2$:
\[
  2\mu\frac{\partial u_r}{\partial r}\bigg|_{x=1}
  = 2\rho\nu\cdot\sigma\cdot\epsilon_0 e^{-\sigma t}\,Y_l^m\,
    \left.\frac{\dd}{\dd x}\!\left(\frac{U}{x^2}\right)\right|_{x=1}
  = 2\rho\nu\sigma\,\epsilon\,Y_l^m\left[\,U'(1) - 2U(1)\,\right].
\]
Dividing through by $\rho\epsilon\,Y_l^m$ and using $\mu = \rho\nu$:
\begin{equation}
  (l-1)(l+2)\frac{T_1}{\rho R} = P_0 - 2\nu\sigma\!\left[\,U'(1) - 2U(1)\,\right].
  \label{eq:BC3}
\end{equation}
\end{derivbox}

\begin{notebox}
BC3 is the only equation containing $T_1$ explicitly. Reid's key move (next section)
is to rescale parameters so that $T_1$ drops out entirely.
\end{notebox}

%% ════════════════════════════════════════════════════════════════════════════
\section{The Characteristic Equation}
\label{sec:char_eq}
%% ════════════════════════════════════════════════════════════════════════════

\subsection{Eliminating $T_1$ via rescaling}

Define the dimensionless ratio
\begin{equation}
  \alpha^2 \equiv \frac{\sigma_{l;0}\,R^2}{\nu},
  \label{eq:alpha_def}
\end{equation}
which measures the inviscid frequency in units of the viscous diffusion rate $\nu/R^2$.
Equivalently, $\alpha = \mathrm{Oh}^{-1/2}\cdot[l(l-1)(l+2)]^{1/4}$ in terms of the
Ohnesorge number.

\begin{derivbox}[Eliminating $T_1$]
From \eqref{eq:inviscid_freq}: $\sigma_{l;0}^2 = l(l-1)(l+2)T_1/(\rho R^3)$,
so $(l-1)(l+2)T_1/(\rho R) = \sigma_{l;0}^2 R^2/l$. Using \eqref{eq:Pi_form},
$P_0 = \sigma^2 R^2\Pi_0/l$. Substituting into BC3 \eqref{eq:BC3}:
\[
  \frac{\sigma_{l;0}^2 R^2}{l}
  = \frac{\sigma^2 R^2\Pi_0}{l} - 2\nu\sigma\!\left[U'(1)-2U(1)\right].
\]
Multiply through by $lR^2/\nu^2$:
\begin{align*}
  \frac{\sigma_{l;0}^2 R^4}{\nu^2}
  &= \frac{\sigma^2 R^4\Pi_0}{\nu^2} - \frac{2l\sigma R^2}{\nu}\!\left[U'(1)-2U(1)\right] \\
  \alpha^4 &= q^4\Pi_0 - 2l\,q^2\!\left[U'(1)-2U(1)\right].
\end{align*}
This can be written compactly as
\begin{equation}
  \alpha^4 = q^2\!\left\{\,q^2\Pi_0 - 2l\!\left[\,U'(1) - 2U(1)\,\right]\right\}.
  \label{eq:alpha4}
\end{equation}
$T_1$ and $\rho$ have cancelled. The problem in terms
of $(\alpha, q)$ is universal.
\end{derivbox}

\subsection{Applying BCs to the general solution}

We now apply BC1 \eqref{eq:BC1} and BC2 \eqref{eq:BC2} to the general solution
\eqref{eq:U_general} to express $C$ and $\Pi_0$ in terms of Bessel functions, then
substitute into \eqref{eq:alpha4} to obtain the characteristic equation.

\paragraph{Applying BC1.}
\begin{equation}
  U(1) = C\,j_l(q) + \Pi_0 = -1.
  \label{eq:bc1_applied}
\end{equation}

\paragraph{Applying BC2.}
Using the results from the note in Section~6.2:
\begin{equation}
  C\!\left[-q^2 j_l + 2(l^2+l-1)j_l - 2q j_l'\right] + 2(l^2-1)\Pi_0 = 0.
  \label{eq:bc2_applied}
\end{equation}

\paragraph{Solving for $C$ and $\Pi_0$.}
From \eqref{eq:bc1_applied}: $\Pi_0 = -1 - Cj_l(q)$. Substituting into
\eqref{eq:bc2_applied} and simplifying:
\[
  C\!\left[(2l - q^2)j_l - 2q j_l'\right] = 2(l^2-1).
\]
Using $q j_l'/j_l = l - qQ_{l+1/2}$ from \eqref{eq:bessel_ratio} and the definition
\eqref{eq:Q_def}:
\begin{equation}
  C = \frac{2(l-1)(l+1)}{(2l-q^2)j_l(q) - 2q\,j_l'(q)}.
  \label{eq:C_soln}
\end{equation}
The expression for $\Pi_0$ follows from \eqref{eq:bc1_applied}.

\paragraph{Substituting into \eqref{eq:alpha4}.}
After substituting the expressions for $C$, $\Pi_0$, $U'(1)$, and $U(1)$ into
\eqref{eq:alpha4}, and expressing all Bessel function ratios in terms of
$Q_{l+1/2}(q) = j_{l+1}(q)/j_l(q)$ via the recurrence \eqref{eq:bessel_ratio},
one obtains the characteristic equation:

\begin{keybox}[The Characteristic Equation (Reid 1960, Eq.~19)]
\begin{equation}
  \frac{\alpha^4}{q^4} + 1
  = \frac{2(l-1)}{q^2}
    \left[\,l + (l+1)\,\frac{q - 2l\,Q_{l+1/2}(q)}{q - 2Q_{l+1/2}(q)}\,\right],
  \label{eq:char_eq}
\end{equation}
where $q^2 = \sigma R^2/\nu$, $\alpha^2 = \sigma_{l;0}R^2/\nu$, and
$Q_{l+1/2}(q) = J_{l+3/2}(q)/J_{l+1/2}(q) = j_{l+1}(q)/j_l(q)$.
\end{keybox}

\begin{notebox}[Universality]
This is identical to Chandrasekhar's characteristic equation for a viscous
\emph{self-gravitating} liquid globe, with the identification that the self-gravitational
parameter maps to $\alpha^2$. The physical restoring force (surface tension vs.\
self-gravity) enters only through $\sigma_{l;0}$, which defines $\alpha$.
\eqref{eq:char_eq} holds for any spherically symmetric restoring force.
\end{notebox}

\begin{notebox}[Algebraic derivation remark]
The algebraic steps from \eqref{eq:C_soln} to \eqref{eq:char_eq} involve expressing
$q^2\Pi_0 - 2l[U'(1)-2U(1)]$ entirely in terms of Bessel ratios. Key simplifications
use the identity
\[
  j_{l+1}(q) = (2l+1)j_l(q)/q - j_{l-1}(q)
\]
(spherical Bessel recurrence) and the relation \eqref{eq:bessel_ratio}.
The algebra is lengthy; each step uses only linear algebra and
these Bessel identities.
\end{notebox}

%% ════════════════════════════════════════════════════════════════════════════
\section{Structure of the Solutions}
%% ════════════════════════════════════════════════════════════════════════════

\subsection{Roots of the characteristic equation}

Equation~\eqref{eq:char_eq} is a transcendental equation for $q^2 = \sigma R^2/\nu$
(or equivalently for $\sigma$), given the physical parameters encoded in $\alpha$.
It has infinitely many solutions.

\paragraph{Physical picture.} The Bessel function $J_{l+1/2}(q)$ in the denominator of
$Q_{l+1/2}$ has infinitely many real zeros $q_1 < q_2 < \cdots$ on the positive real axis.
Near each zero the characteristic equation \eqref{eq:char_eq} has a pole, and each pole
``captures'' one pair of roots. These correspond to increasingly high radial overtones of
the oscillation, i.e., modes with more and more nodal surfaces in the radial direction
inside the drop.

\paragraph{Ordering by decay rate.} Roots with \emph{smaller} real part of $\sigma$ (i.e.,
slower decay) are the physically dominant ones: they are the last to die out and govern
the observable oscillation of the drop. The two roots with the smallest $\mathrm{Re}(\sigma)$
correspond to the fundamental surface oscillation at harmonic order $l$; all higher roots
decay exponentially faster.

\begin{notebox}[Sign convention caution]
In Reid's convention, $\epsilon\sim e^{-\sigma t}$ with $\mathrm{Re}(\sigma)>0$
meaning decay. So the \emph{smallest positive} real part corresponds to the
\emph{slowest} decay, i.e., the dominant mode. This is the opposite of the
situation in stability analysis where one looks for the most-positive growth rate.
\end{notebox}

\subsection{Two regimes: oscillatory vs.\ aperiodic}

For fixed $l$, the character of the two dominant roots depends on $\alpha^2$:

\begin{itemize}
  \item \textbf{Large $\alpha^2$ (low viscosity, large drop):} the two dominant roots
        are complex conjugates, $\sigma = \gamma \pm \ii\omega_d$, corresponding to
        \emph{damped oscillations}.
  \item \textbf{Small $\alpha^2$ (high viscosity, small drop):} the two dominant roots
        are both real and positive, corresponding to two \emph{aperiodic decaying modes}
        (overdamped).
\end{itemize}

The transition occurs at a critical value of $\alpha^2$. For $l=2$, this critical point
was determined numerically by Chandrasekhar:
\begin{equation}
  \sigma_{2;0}R^2/\nu = 3.69, \qquad \sigma_{2;\nu}/\sigma_{2;0} = 0.968.
  \label{eq:critical}
\end{equation}
For water in air ($T_1 = 74\;\mathrm{dyn/cm}$, $\nu = 0.01\;\mathrm{cm^2/s}$), this
gives a critical radius $R_c \approx 0.23\;\mathrm{mm}$.

\begin{itemize}
  \item Drops with $R > R_c$: oscillate with damping.
  \item Drops with $R < R_c$: aperiodically damped (no oscillations).
\end{itemize}

\subsection{The small-viscosity limit}

As $\nu \to 0$ (or $q \to \infty$), the spherical Bessel functions have the large-argument
asymptotics $J_\nu(q) \sim \sqrt{2/(\pi q)}\cos(q - \nu\pi/2 - \pi/4)$, which for
$\nu = l+\tfrac{1}{2}$ reduces to $J_{l+1/2}(q) \sim \sqrt{2/(\pi q)}\sin(q - l\pi/2)$.
The ratio $Q_{l+1/2}(q) = j_{l+1}(q)/j_l(q)$ therefore oscillates as $q\to\infty$
(with poles at every zero of $j_l$), but $Q_{l+1/2}(q)/q \to 0$ between poles,
since the amplitude of $J_{l+3/2}/J_{l+1/2}$ remains $O(1)$ away from zeros.
The bracket on the right-hand side of \eqref{eq:char_eq},
$l + (l+1)(q - 2l\,Q_{l+1/2})/(q - 2Q_{l+1/2})$, approaches $l + (l+1)\cdot 1 = 2l+1$
as $q\to\infty$. The characteristic equation \eqref{eq:char_eq} then simplifies to:
\[
  \frac{\alpha^4}{q^4} + 1 = \frac{2(l-1)(2l+1)}{q^2},
\]
which is a quadratic in $q^2/\alpha^2$, yielding
\begin{equation}
  \sigma_{l;\nu} = (l-1)(2l+1)\,\frac{\nu}{R^2} \pm \ii\,\sigma_{l;0},
  \label{eq:lamb}
\end{equation}
recovering \emph{Lamb's classical result}: viscosity adds a real damping rate
$(l-1)(2l+1)\nu/R^2$ to the inviscid oscillation frequency $\pm\ii\sigma_{l;0}$.

%% ════════════════════════════════════════════════════════════════════════════
\section{Connection to Mola\v{c}ek \& Bush (2012)}
%% ════════════════════════════════════════════════════════════════════════════

In their quasi-static model of drop impact, Mola\v{c}ek \& Bush parameterize the
kinetic energy and viscous dissipation associated with each surface harmonic mode $m$
(their notation) via two coefficients $A_m(\mathrm{Oh}_m)$ and $D_m(\mathrm{Oh}_m)$,
which appear in (their Eq.~31):
\[
  \mathrm{K.E.} = \pi\rho R_0^5 \dot{B}^2 \sum_m A_m \frac{2b_m^2}{m(2m+1)},
  \qquad
  D = 8\pi\mu R_0^3 \dot{B}^2 \sum_m D_m \frac{m}{2m+1}b_m^2.
\]

These coefficients are defined so that the two roots of the quadratic
\begin{equation}
  A_m b^2 - 2a D_m b + 1 = 0, \qquad a \equiv \mathrm{Oh}\sqrt{m(m-1)(m+2)},
\end{equation}
reproduce the two dominant roots of Reid's characteristic equation \eqref{eq:char_eq}
for harmonic order $l=m$. The variable identification connecting MB's quadratic to Reid's
transcendental equation is $q^2 = (b/a)\cdot\alpha^2$ (i.e.\ $b = a\cdot q^2/\alpha^2$),
under which MB's $W(b/a) = Q_{l+1/2}(q)/q$. In other words, $A_m$ and $D_m$ encode the
\emph{result} of solving Reid's problem at each order $m$, compressed into two numbers
per mode via a quadratic fit to the dominant eigenvalue pair.

\paragraph{Workflow.}
\begin{enumerate}
  \item For each mode $m$ and Ohnesorge number $\mathrm{Oh}$, evaluate
        $\alpha^2 = \mathrm{Oh}^{-1}\sqrt{m(m-1)(m+2)}$ and $q^2 = \sigma R^2/\nu$.
  \item Numerically solve \eqref{eq:char_eq} for the two roots $b_{1,2}$ with the
        smallest $\mathrm{Re}(\sigma)$.
  \item Read off $A_m$ and $D_m$ by matching the quadratic:
        sum of roots $= 2aD_m/A_m$; product of roots $= 1/A_m$.
  \item Tabulate or fit $A_m(\mathrm{Oh}_m)$ and $D_m(\mathrm{Oh}_m)$ for use in
        the Lagrangian equation of motion.
\end{enumerate}

\paragraph{Justification for discarding higher roots.} The quasi-static assumption
restricts the drop shape to evolve through a one-parameter family of equilibrium shapes.
Within that family, only the fundamental surface mode at each harmonic order $l$ is
representable. Higher radial overtones have internal nodal surfaces absent from the quasi-static shape family. They decay on timescales $\ll (R^2/\nu)$, well separated from the impact timescale.

%% ════════════════════════════════════════════════════════════════════════════
\section{Summary of Key Equations}
%% ════════════════════════════════════════════════════════════════════════════

\begin{center}
\begin{tabular}{ll}
  \toprule
  \textbf{Equation} & \textbf{Physical content} \\
  \midrule
  $\sigma_{l;0}^2 = l(l-1)(l+2)\,T_1/(\rho R^3)$ & Inviscid capillary frequency \\
  $q^2 = \sigma R^2/\nu$ & Complex decay rate, scaled \\
  $\alpha^2 = \sigma_{l;0}R^2/\nu$ & Inviscid frequency, scaled \\
  $Q_{l+1/2}(q) = J_{l+3/2}(q)/J_{l+1/2}(q)$ & Bessel function ratio \\
  $\alpha^4/q^4 + 1 = (2(l-1)/q^2)[\cdots]$ & Reid's characteristic equation \\
  $\sigma = (l-1)(2l+1)\nu/R^2 \pm \ii\,\sigma_{l;0}$ & Small-$\nu$ (Lamb) limit \\
  \bottomrule
\end{tabular}
\end{center}

\bigskip

%% ════════════════════════════════════════════════════════════════════════════
\section{Weakly Nonlinear Shear-Thinning Drops: the Carreau Model}
\label{sec:carreau}
%% ════════════════════════════════════════════════════════════════════════════

Sections 1--9 treat a Newtonian fluid; Section~\ref{sec:oldroydB} treats an
Oldroyd-B fluid, in both cases at linear order in the oscillation amplitude
$\epsilon$. This section derives the leading-order effect of shear-thinning
using the Carreau constitutive law. Shear-thinning is a nonlinear effect: the
viscosity depends on the instantaneous strain rate. It does not enter the
linear eigenvalue problem, but first appears at $O(\epsilon^3)$ as a cubic
correction to the mode damping.

%% ────────────────────────────────────────────────────────────────────────────
\subsection{The Carreau constitutive law}
%% ────────────────────────────────────────────────────────────────────────────

The Carreau model relates effective viscosity to the scalar shear rate
$\dot\gamma = \sqrt{2\,e_{ij}e_{ij}}$:
\begin{equation}
  \mu_{\mathrm{eff}}(\dot\gamma)
  = \mu_0 \left[1 + (\lambda_c\,\dot\gamma)^2\right]^{(n-1)/2},
  \label{eq:carreau}
\end{equation}
where $\mu_0$ is the zero-shear-rate viscosity, $\lambda_c$ is the Carreau
relaxation time, and $n \in (0,1]$ is the power-law index. Setting $n=1$
gives $\mu_{\mathrm{eff}} = \mu_0$ (Newtonian); for $n < 1$ the viscosity
decreases with shear rate (shear-thinning). The shear-thinning parameter
\begin{equation}
  \varepsilon_{ST} \equiv \frac{1-n}{2} \geq 0
  \label{eq:epsST}
\end{equation}
is zero for a Newtonian fluid and positive for $n < 1$.

%% ────────────────────────────────────────────────────────────────────────────
\subsection{Perturbation structure}
%% ────────────────────────────────────────────────────────────────────────────

In the linearised flow, $\mathbf{e} = O(\epsilon)$ and hence
$\dot\gamma = O(\epsilon)$. Expanding \eqref{eq:carreau} in $\epsilon$:
\begin{equation}
  \mu_{\mathrm{eff}}
  = \mu_0\left[1 - \varepsilon_{ST}(\lambda_c\dot\gamma)^2 + O(\epsilon^4)\right].
  \label{eq:carreau_expand}
\end{equation}
The $O(\epsilon)$ viscosity correction is exactly zero. The first non-trivial
Carreau correction appears at $O(\epsilon^2)$ in $\mu_{\mathrm{eff}}$. Because
the stress is $\boldsymbol\tau = 2\mu_{\mathrm{eff}}\,\mathbf{e}$ and
$\mathbf{e} = O(\epsilon)$, the stress correction enters at $O(\epsilon^3)$.
The linear problem is therefore identical to the Newtonian problem with
viscosity $\mu_0$; the eigenvalue, eigenfunction, and characteristic equation
of Sections 1--9 are unchanged.

%% ────────────────────────────────────────────────────────────────────────────
\subsection{Shear-rate field and the correction integral $\Gamma_l$}
%% ────────────────────────────────────────────────────────────────────────────

At $O(\epsilon)$ the velocity field is the Newtonian one. For a freely
oscillating drop in mode $l$, the radial and tangential components are:
\begin{equation}
  u_r = \dot{b}_l\,\frac{U(x)}{x^2}\,P_l(\cos\theta),
  \qquad
  u_\theta = \dot{b}_l\,\frac{f(x)}{x^2}\,\frac{dP_l}{d\theta},
  \label{eq:vel_carreau}
\end{equation}
where $b_l$ is the mode amplitude and $f(x)$ follows from incompressibility.

Define the correction integral as the volume and time average of
$\dot\gamma^4$, normalized by $(\dot{b}_l/R)^2$ and the mode norm
$\mathcal{N}_l = \int_0^1 |U(x)|^2 x^2\,dx$:
\begin{equation}
  \Gamma_l \equiv
  \frac{1}{\mathcal{N}_l}
  \int_0^1
  \left\langle\int_0^\pi \left(\frac{\dot\gamma(x,\theta)}{\dot{b}_l/R}\right)^4
  \sin\theta\,d\theta\right\rangle_t
  x^2\,dx,
  \label{eq:Gamma_def}
\end{equation}
where $\langle\cdot\rangle_t$ denotes the time average over one period. In the
inviscid limit ($\mathrm{Oh}\to 0$, real $q$), this integral is computable
analytically from the Newtonian eigenfunctions. For $l=2$:
\begin{equation}
  \Gamma_2 = \frac{1783566}{385} \approx 4632.6
  \qquad (\mathrm{Oh}\to 0).
  \label{eq:Gamma2}
\end{equation}
At finite $\mathrm{Oh}$, $q$ is complex and $U(x)$ is complex-valued. The
time-averaged $\langle\dot\gamma^4\rangle_t$ then picks up cross-terms between
the real and imaginary parts of the eigenfunction. The general finite-$\mathrm{Oh}$
expression is derived in the companion notebook
(\texttt{shear\_thinning\_derivation.ipynb}).

%% ────────────────────────────────────────────────────────────────────────────
\subsection{Modified mode equation}
%% ────────────────────────────────────────────────────────────────────────────

The Carreau correction adds extra viscous dissipation through the Rayleigh
dissipation function:
\[
  \delta\Phi = \mu_0\,\varepsilon_{ST}\,\lambda_c^2
  \int_V \dot\gamma^4\,dV.
\]
The generalized force $\partial(\delta\Phi)/\partial\dot{b}_l$ is cubic in
$\dot{b}_l$, and moves to the right-hand side of the mode equation with a
positive sign (shear-thinning reduces dissipation):
\begin{equation}
  A_l\,\ddot{b}_l
  + D_l^{(0)}\,\dot{b}_l
  + \omega_l^2\,b_l
  = \varepsilon_{ST}\,\Lambda^2\,\Gamma_l\,D_l^{(0)}\,|\dot{b}_l|^2\,\dot{b}_l
  + O(\epsilon^5),
  \label{eq:mode_carreau}
\end{equation}
where
\begin{align}
  D_l^{(0)} &= 2\,\mathrm{Oh}\,(l-1)(2l+1), \label{eq:Dl0}\\
  \Lambda    &= \lambda_c\,\sigma_{l;0}, \label{eq:Lambda}\\
  \omega_l^2 &= l(l-1)(l+2). \notag
\end{align}
Here $D_l^{(0)}$ is the Newtonian damping coefficient (from the Lamb limit),
$\Lambda$ is the Carreau number (Carreau relaxation time non-dimensionalised
by the inviscid period), and $A_l$ is the kinetic-energy coefficient from
Mola\v{c}ek \& Bush~\cite{molacek2012}.

Moving the cubic term to the left defines an amplitude-dependent effective
damping:
\begin{equation}
  D_{\mathrm{eff}} = D_l^{(0)}
  \left[1 - \varepsilon_{ST}\,\Lambda^2\,\Gamma_l\,|\dot{b}_l|^2\right].
  \label{eq:Deff}
\end{equation}
For $\varepsilon_{ST} > 0$, $D_{\mathrm{eff}} < D_l^{(0)}$ at finite
amplitude, so the drop damps more slowly than a Newtonian drop of the same
zero-shear viscosity.

%% ────────────────────────────────────────────────────────────────────────────
\subsection{Slow amplitude equation}
\label{sec:carreau_slow}
%% ────────────────────────────────────────────────────────────────────────────

For $\mathrm{Oh} \ll 1$, the leading-order solution is
$b_l \approx a(t)\cos(\omega_l t + \phi)$ with slowly varying amplitude $a(t)$.
The cubic term in \eqref{eq:mode_carreau} is a resonant forcing at the
oscillation frequency. Applying the method of averaging (projecting onto
$\sin\omega_l t$ and integrating over one period) gives:

\begin{keybox}[Slow amplitude equation (Carreau drop)]
\begin{equation}
  \frac{da}{dt}
  = -\gamma_l^{(0)}\,a
  \left[1 - \frac{3}{4}\,\varepsilon_{ST}\,\Lambda^2\,\Gamma_l\,a^2\right],
  \label{eq:slow_amp}
\end{equation}
where $\gamma_l^{(0)} = (l-1)(2l+1)\,\mathrm{Oh}$ is the Newtonian decay rate.
\end{keybox}

The factor $3/4$ is the time average of $\cos^2(\omega_l t)\sin^2(\omega_l t)$
over one period. Equation~\eqref{eq:slow_amp} is valid to
$O(\varepsilon_{ST}\Lambda^2 a^2)$. The bracket is less than unity for
$\varepsilon_{ST} > 0$, so the instantaneous decay rate $|da/dt|/a$ is
reduced at every amplitude relative to the Newtonian value $\gamma_l^{(0)}$.
There is no non-zero fixed point: $a\to 0$ as $t\to\infty$, but more slowly
than in the Newtonian case.

\begin{notebox}[Scope]
Equation~\eqref{eq:slow_amp} requires:
\begin{itemize}
  \item $\mathrm{Oh} \ll 1$: lightly damped oscillations, so the method of
        averaging applies.
  \item $\varepsilon_{ST}\Lambda^2 a^2 \ll 1$: the bracket in
        \eqref{eq:slow_amp} must remain positive.
  \item No contact: $\Gamma_l$ is computed from the free-oscillation Reid
        velocity field. Near a contact line $\dot\gamma\to\infty$ and the
        expansion $(\lambda_c\dot\gamma)^2\ll 1$ breaks down.
\end{itemize}
The Julia solver (\texttt{DropSolver.jl}) implements Carreau rheology via
\texttt{STParams}: pre-compute $\Gamma_l$ per mode from
\texttt{notebooks/shear\_thinning\_derivation.ipynb} and pass it to
\texttt{solve\_drop!} as a keyword argument.
\end{notebox}


%% ════════════════════════════════════════════════════════════════════════════
\section{General-Exponent Carreau-Yasuda: Why Amplitude Perturbation Theory
         Fails, and What Survives}
\label{sec:cy_firstprinciples}
%% ════════════════════════════════════════════════════════════════════════════

Section~\ref{sec:carreau} derives a slow amplitude equation for the classical
Carreau law, $\mu_{\mathrm{eff}}=\mu_0[1+(\lambda_c\dot\gamma)^2]^{(n-1)/2}$,
valid for $\mathrm{Oh}\ll 1$ and small oscillation amplitude. That derivation
is correct for the model it treats. This section asks a harder question:
does an analogous small-amplitude expansion exist for the \emph{general}
Carreau-Yasuda law,
\begin{equation}
  \eta(\dot\gamma) = \eta_\infty + (\eta_0-\eta_\infty)\bigl[1+(\lambda_c\dot\gamma)^a\bigr]^{(n-1)/a},
  \label{eq:cy_general}
\end{equation}
with the shape exponent $a$ left general rather than fixed at $2$ -- the form
this repo's own fitted validation fluid actually uses
($\mathrm{Oh}_0\approx 57.37$, $\lambda_c\approx 3.05\times 10^4$,
$a\approx 0.743$, $\varepsilon_{ST}\equiv(1-n)/a\approx 0.9996$)?

The short answer, reached only after an extensive derivation and two
independent adversarial checks of every step, is: \textbf{no, not for this
fluid's fitted parameters}. This is not a failure of nerve or a shortcut --
it is a precisely quantified mathematical fact about where this particular
fluid's physics actually lives on the Carreau-Yasuda curve. What survives
the investigation is three genuinely new, general, CAS-verified results that
have nothing to do with any specific fluid's parameters, plus a clear,
quantitative diagnosis of exactly what a first-principles treatment of
\emph{this} fluid requires instead.

%% ────────────────────────────────────────────────────────────────────────────
\subsection{Why the classical $O(\epsilon^3)$ route is not available here}
\label{sec:cy_why_not_epsilon3}
%% ────────────────────────────────────────────────────────────────────────────

Section~\ref{sec:carreau}'s expansion rests on one fact: with $a=2$ fixed,
$(\lambda_c\dot\gamma)^2$ is an even integer power of $\dot\gamma$, hence an
analytic (Taylor-expandable) function of the oscillation amplitude
$\epsilon$, since $\dot\gamma=O(\epsilon)$ linearly in Reid's single-mode
ansatz. The viscosity correction is then $O(\epsilon^2)$, the stress
correction $O(\epsilon^3)$, and a standard Landau-equation-style
method-of-averaging calculation goes through cleanly.

\begin{notebox}[The exponent matters, not just its size]
For a \emph{general} exponent $a$, $(\lambda_c\dot\gamma)^a\sim
(\lambda_c\epsilon)^a\,\hat g(\theta)^a$ for some $O(1)$ angular shape
$\hat g$. If $a$ is not an even integer, $|\epsilon|^a$ is not differentiable
at $\epsilon=0$ (it is only H\"older continuous there). For this fluid's
fitted $a\approx 0.743 <1$, the viscosity correction scales as
$|\epsilon|^{0.743}$, which is \emph{larger} than $\epsilon$ itself for small
$\epsilon$ (since $|\epsilon|^{a-1}\to\infty$ as $\epsilon\to0$ when $a<1$).
There is no regime, however small the amplitude, in which this correction is
a small perturbation of a fixed integer-power hierarchy: it sits at a
genuinely fractional order between Reid's $O(\epsilon)$ linear terms and the
usual $O(\epsilon^2)$ geometric mode-coupling correction.
\end{notebox}

This alone does not doom a perturbative treatment -- a fractional-order
expansion parameter is unusual but not unheard of. What actually closes off
the small-amplitude route for this fluid is quantitative, not structural,
and is established in the next section.

%% ────────────────────────────────────────────────────────────────────────────
\subsection{Two natural anchors, and why neither supports a small correction}
\label{sec:cy_anchors}
%% ────────────────────────────────────────────────────────────────────────────

A shear-thinning fluid's viscosity interpolates between two Newtonian
limits: $\eta_0$ at rest ($\dot\gamma\to0$) and $\eta_\infty$ at infinite
shear ($\dot\gamma\to\infty$). Reid's characteristic equation
(Section~\ref{sec:char_eq}) gives the \emph{exact} damping rate
$\lambda_l(\mathrm{Oh})$ and squared frequency $\omega_l^2(\mathrm{Oh})$ at
\emph{any} Ohnesorge number, so either limit is a legitimate place to try to
anchor a correction.

\begin{derivbox}[Anchoring at rest, $\mathrm{Oh}_0$: fails outright]
At $\mathrm{Oh}_0\approx 57.37$, every mode this fluid's solver resolves is
heavily overdamped:
\[
  \begin{array}{cccc}
    l & \lambda_l & \omega_l^2 & \\
    2 & 203.8 & 7.48 & \text{overdamped} \\
    3 & 459.4 & 25.48 & \text{overdamped} \\
    5 & 1119.3 & 102.9 & \text{overdamped} \\
    10 & 3681.4 & 665.4 & \text{overdamped}
  \end{array}
\]
There is no oscillation to average over, so a multiple-scales/envelope
treatment has no fast carrier wave to modulate. Separately, solving for the
shear rate at which the viscosity correction reaches even the $1\%$ level
gives $\dot\gamma\lesssim 7\times10^{-8}$ -- a shear rate no resolved mode
velocity in this solver ever produces. Both the qualitative premise (an
oscillator to perturb) and the quantitative premise (a small correction)
fail simultaneously.
\end{derivbox}

\begin{derivbox}[Anchoring at infinite shear, $\mathrm{Oh}_\infty$: better
                 posed, but still not small]
Define $\mathrm{Oh}_\infty\equiv\mathrm{Oh}_0\,(\eta_\infty/\eta_0) =
\mathrm{Oh}_0(1-\varepsilon_{ST})\approx 0.0254$. Unlike $\mathrm{Oh}_0$,
this \emph{is} comfortably underdamped: at $l=2$,
$\lambda_2=0.116$, $\omega_2^2=7.95$. Reid's oscillator convention is
$\ddot A+2\lambda\dot A+\omega^2A=\text{forcing}$
(\texttt{julia/src/reid.jl}), so the quality factor is
$Q\equiv\sqrt{\omega_2^2}/(2\lambda_2)\approx 12$ -- a genuine separation
between the decay time and the oscillation period, exactly the condition a
weakly-damped multiple-scales treatment needs. A live solver trace (three
impact energies, fine time resolution) confirms the drop's actual effective
Ohnesorge number, once excited, is permanently pinned away from
$\mathrm{Oh}_0$ after a single clean transient at first contact, with no
recurrence on later contact/lift-off cycles.

But ``pinned away from $\mathrm{Oh}_0$'' is not the same claim as ``close to
$\mathrm{Oh}_\infty$.'' Evaluating the exact Carreau-Yasuda law at the shear
rates this fluid's live trace actually visits ($\dot\gamma\sim
0.05$--$1.5$, in the solver's nondimensional units) gives
$\mathrm{Oh}_{\mathrm{eff}}/\mathrm{Oh}_\infty$ ratios of roughly $2$ to
$7$ -- not a small departure from the anchor. And the relevant question is
not this ratio directly, but whether the \emph{first-order Taylor expansion}
of Reid's own exact $\lambda_l(\mathrm{Oh})$, $\omega_l^2(\mathrm{Oh})$
around $\mathrm{Oh}_\infty$ reproduces their true values over that range.
Checked directly against the exact (continuation-solved) relations, at four
points spanning the operative band, $\mathrm{Oh}=0.03,0.05,0.08,0.1$
(ratios $1.2$--$3.9\times\mathrm{Oh}_\infty$): $\omega_l^2$ is nearly flat in
$\mathrm{Oh}$ near $\mathrm{Oh}_\infty$ -- linear-approximation error stays
under $0.3\%$ for $l=2$ across the whole band, and is still small for
higher $l$, growing to about $1.5\%$ by $l=10$ at the band's upper end. But
$\lambda_l$ is far more sensitive and its error \emph{grows across the
band}, not just with $l$: at the low end ($\mathrm{Oh}=0.03$) every $l$
tested is under $0.3\%$, but by the upper end ($\mathrm{Oh}=0.1$) the error
has grown to $8\%$ ($l=2$), $12\%$ ($l=3$), $17\%$ ($l=5$), and $26\%$
($l=10$) -- the same growing-with-$l$ curvature that broke the
$\mathrm{Oh}_0$ anchor reappears here, mode-dependently, and here it also
grows across the operative band itself, not just across modes.
Moving the anchor to some other, empirically ``typical'' $\mathrm{Oh}^\ast$
does not rescue this: the curvature in $\lambda_l(\mathrm{Oh})$ at high $l$
sets in almost immediately away from any point in the operative range, so no
placement of a single anchor buys a genuinely wide small-parameter regime.
\end{derivbox}

\begin{keybox}[The honest conclusion]
For this fluid's fitted parameters, there is no single reference viscosity
-- rest, infinite-shear, or any point in between -- around which a
regular (or matched-asymptotic, single-anchor) perturbation expansion of
Reid's damping rate is uniformly valid across the modes this solver
resolves. The physics genuinely lives in the fully nonlinear, transitional
part of the Carreau-Yasuda curve, not in a neighborhood of either plateau.
This is a quantitative fact about \emph{this fluid's} $\lambda_c$ and $a$,
not a universal statement about shear-thinning drops.
\end{keybox}

%% ────────────────────────────────────────────────────────────────────────────
\subsection{What Reid's problem already gives you for free}
\label{sec:cy_exact_evaluation}
%% ────────────────────────────────────────────────────────────────────────────

Because the failure above is specifically a failure of \emph{linearizing}
$\lambda_l(\mathrm{Oh})$, $\omega_l^2(\mathrm{Oh})$, and because these
functions are already known \emph{exactly}, for any $\mathrm{Oh}$, via
Reid's characteristic equation (solved numerically by continuation, and
tabulated for fast repeated evaluation), the correct response is not to
build a better linear correction -- it is to not linearize this piece at
all. Given an instantaneous effective Ohnesorge number
$\mathrm{Oh}_{\mathrm{eff}}(t)$ from any closure (including the one derived
in Section~\ref{sec:cy_temporal} below), the exact, non-perturbative
evaluation of $\lambda_l(\mathrm{Oh}_{\mathrm{eff}}(t))$,
$\omega_l^2(\mathrm{Oh}_{\mathrm{eff}}(t))$ has no validity-range caveat and
is already the practice this repo's existing (heuristic) Carreau-Yasuda
extension follows. That part of the existing scheme was never the weak
point. The weak point is upstream of it: whether it is valid to replace a
velocity field's genuinely \emph{spatially varying} local viscosity with a
single scalar $\mathrm{Oh}_{\mathrm{eff}}(t)$ fed into a formula built for a
spatially \emph{uniform} viscosity. That is the question the rest of this
section answers.

%% ────────────────────────────────────────────────────────────────────────────
\subsection{Machinery I: adjoint sensitivity for Reid's boundary-value problem}
\label{sec:cy_adjoint}
%% ────────────────────────────────────────────────────────────────────────────

Reid's velocity equation, $\mathcal{L}[U]\equiv U''-l(l+1)U/x^2+q^2U$
(Section~\ref{sec:velocity_ode}), is formally self-adjoint on $(0,1)$ with
the plain $L^2$ inner product -- no weight function is needed, because
$\mathcal{L}$ has no first-derivative term:
\begin{derivbox}[Lagrange identity for $\mathcal{L}$]
For any twice-differentiable $U,W$,
\begin{equation}
  W\,\mathcal{L}[U] - U\,\mathcal{L}[W]
  = \frac{d}{dx}\Bigl[\,W\,U' - W'\,U\,\Bigr].
  \label{eq:lagrange}
\end{equation}
Verified symbolically (Symbolics.jl, abstract $U(x)$, $W(x)$): both sides
agree identically.
\end{derivbox}

This gives a shortcut for exactly the calculation a perturbation of Reid's
problem needs: the boundary derivative of a particular solution, without
solving the ODE.

\begin{keybox}[Adjoint shortcut]
Let $Y(x)$ be regular at $x=0$, satisfy $Y(1)=0$, and solve
$\mathcal{L}[Y]=\mathrm{RHS}(x)$ for a known forcing $\mathrm{RHS}$, at a
fixed wavenumber $q_0$ (so $\mathcal{L}$ uses $q_0^2$). Then
\begin{equation}
  Y'(1) = \frac{1}{j_l(q_0)}\int_0^1 x\,j_l(q_0x)\,\mathrm{RHS}(x)\,dx,
  \label{eq:adjoint_shortcut}
\end{equation}
using the regular homogeneous solution $x\,j_l(q_0x)$ as the adjoint test
function (both boundary terms at $x=0$ vanish by regularity; the boundary
term at $x=1$ collapses using $Y(1)=0$).
\end{keybox}

\begin{notebox}[Numerically verified, not just derived]
Checked against direct numerical shooting (regular-solution seeding near
$x=0$, RK4 integration, homogeneous-solution correction to enforce
$Y(1)=0$) for three independent $(l,q_0,\mathrm{RHS})$ combinations: relative
agreement of $10^{-9}$ to $10^{-11}$ once the shooting method's own
near-origin seed is chosen at a moderate radius ($x_0\sim 10^{-2}$, not
$10^{-6}$ -- the regular solution scales as $x^{l+1}$ near the origin, and
seeding too close to $0$ loses precision to floating-point underflow before
the adjoint comparison is even meaningful).
\end{notebox}

Reid's boundary conditions close the system: BC1 ($U(1)=-1$) is unaffected
by any viscosity correction, since it is purely kinematic, and BC2
($\mathcal L_2[U](1)=0$ plus a known correction) is linear in the same
unknowns the shortcut already produces. Combining the two turns the search
for the $O(\delta)$ shift in $q^2$, from \emph{any} new forcing added to
Reid's ODE and both stress boundary conditions, into an ordinary linear
solve -- no further ODE integration is required. This reduction has been
worked through by hand for the general case (not yet promoted to an
assertion in the backing script, since the resulting expression is a large
symbolic fraction rather than a compact closed form); what \emph{is}
CAS-verified end to end is the shortcut itself
(Eq.~\ref{eq:adjoint_shortcut}), which is the genuinely nontrivial step.

%% ────────────────────────────────────────────────────────────────────────────
\subsection{Machinery II: the strain-rate field of Reid's actual viscous mode}
\label{sec:cy_strain}
%% ────────────────────────────────────────────────────────────────────────────

Any first-principles shear-rate calculation must be built from Reid's actual
\emph{viscous} velocity profile $U(x)=C\,x\,j_l(qx)+\Pi_0x^{l+1}$
(Section~\ref{sec:velocity_ode}), not the inviscid potential-flow shape
$r^lP_l(\cos\theta)$ that this repo's existing heuristic scripts use to
estimate a characteristic shear rate. The two are not interchangeable: Reid's
own damping normalization comes from the homogeneous (Bessel) part of
$U(x)$, which \emph{is} the viscous correction to potential flow -- using
the potential-flow shape to estimate shear rate is inconsistent with Reid's
theory already at zeroth order, before any shear-thinning correction is
applied.

Writing $u_r=F(x)P_l(\cos\theta)$, $F(x)\equiv U(x)/x^2$, and $u_\theta$
from the stream-function relation of Section~\ref{sec:bc2}
($u_\theta=-(2F+xF')\sin\theta\,P_l'(\cos\theta)/[l(l+1)]$), the standard
axisymmetric strain-rate tensor is
\begin{align}
  e_{rr}   &= F'(x)\,P_l(\mu), \\
  e_{\theta\theta} &= \frac{1}{x}\frac{\partial u_\theta}{\partial\theta} + \frac{u_r}{x}, \\
  e_{\varphi\varphi} &= \frac{u_r + u_\theta\cot\theta}{x}, \\
  e_{r\theta} &= \frac12\left[x\frac{\partial}{\partial x}\!\left(\frac{u_\theta}{x}\right) + \frac{1}{x}\frac{\partial u_r}{\partial\theta}\right],
  \qquad \mu\equiv\cos\theta.
  \label{eq:strain_tensor}
\end{align}

\begin{notebox}[Incompressibility check]
$e_{rr}+e_{\theta\theta}+e_{\varphi\varphi}=0$ identically, for \emph{any}
radial profile $F(x)$ -- verified both symbolically (concrete Legendre
polynomials, $l=2,3$, after resolving a $\sin^2\theta+\cos^2\theta=1$
simplification Symbolics.jl does not apply automatically) and numerically
(concrete polynomial $F$, $l=2,\dots,8$, to floating-point precision at
every tested $(x,\theta)$). This is a kinematic identity of the poloidal
representation itself, confirming the strain-tensor construction is
consistent regardless of what $F(x)$ turns out to be.
\end{notebox}

%% ────────────────────────────────────────────────────────────────────────────
\subsection{Machinery III: the period-$\pi$ lemma}
\label{sec:cy_period_pi}
%% ────────────────────────────────────────────────────────────────────────────

At $\mathrm{Oh}_\infty$, $q$ is genuinely complex ($q_0\approx
7.60-7.30\,\mathrm{i}$ for $l=2$), so Reid's mode is a true damped
oscillation, not a pure decay. The physical (real) strain field at phase
$\phi\equiv\omega t$ is $\mathrm{Re}\bigl[e_{ij}(x,\theta)\,e^{-\mathrm{i}\phi}\bigr]$
for each complex tensor component $e_{ij}$ built from~\eqref{eq:strain_tensor}
with complex $F$. A natural first attempt at a shear-thinning correction --
evaluate everything at one representative phase, e.g. $\phi=0$ -- turns out
to be indefensible: numerically, the real physical shape at $x_0=0.7$
swings from $-0.345$ to $+0.345$ across one period, nowhere near constant.
Any nonlinear (in particular, fractional-power) function of this field
therefore depends on \emph{which} phase was chosen, unless the choice is
justified by an actual time-average.

\begin{keybox}[The shear-rate invariant cannot resonate at the mode's own frequency]
Let $S\equiv\sqrt{2\,e_{ij}e_{ij}}$ (summed over the strain-tensor
components) be built from any single-frequency oscillating axisymmetric
poloidal field, real or complex amplitude. Then $S(\phi)$ is exactly
periodic with period $\pi$ in $\phi=\omega t$, not $2\pi$, and hence its
Fourier series over the full period contains only even harmonics: the
coefficient of $\cos(\phi)$ and $\sin(\phi)$ (the fundamental, $m=1$) is
identically zero, for $S$ itself and for \emph{any} function of $S$
(fractional power or otherwise).
\end{keybox}

\begin{derivbox}[Proof]
Each strain component satisfies
$\mathrm{Re}(e_{ij}e^{-\mathrm{i}\phi})=|e_{ij}|\cos(\phi-\arg e_{ij})$.
Squaring, $|e_{ij}|^2\cos^2(\phi-\arg e_{ij}) =
\tfrac12|e_{ij}|^2\bigl[1+\cos(2\phi-2\arg e_{ij})\bigr]$, which is
manifestly invariant under $\phi\to\phi+\pi$. $S^2$ is a sum of such terms,
hence itself period-$\pi$; $S=\sqrt{S^2}\ge0$ and any real function $h(S)$
inherit the period exactly (no branch ambiguity, since $S^2\ge 0$
throughout). A period-$\pi$ function integrated against $\sin(m\phi)$ or
$\cos(m\phi)$ over $[0,2\pi)$ vanishes identically for every odd $m$, since
the two half-period copies enter with opposite sign and cancel.
\end{derivbox}

\begin{notebox}[Confirmed two ways]
Numerically: $S^2(\phi+\pi)=S^2(\phi)$ to $10^{-15}$--$10^{-16}$ at multiple
independent $(x,\theta)$ points. Via direct Fourier projection on a joint
$800\times800$ $(\theta,\phi)$ grid: the $m=1$ coefficient (both $\cos$ and
$\sin$) is zero to $10^{-17}$--$10^{-18}$ at every spherical-harmonic degree
$l'$ tested ($0,2,4,6,8$), while $m=0$ and $m=2$ are the genuinely nonzero
temporal channels.
\end{notebox}

This has an immediate structural consequence: since no forcing term
resonant at the base mode's own frequency exists, the adjoint machinery of
Section~\ref{sec:cy_adjoint} -- built for exactly this kind of resonant
solvability condition -- has nothing to act on at $m=1$. The leading
temporal effect of any generalized-Newtonian correction on this mode is
therefore carried entirely by the $m=0$ (period-averaged, effectively
steady) channel; the $m=2$ channel represents a distinct, second-harmonic
\emph{parametric} coupling (a $2\!:\!1$ resonance structure with the base
oscillation), a different physical phenomenon not addressed here.

%% ────────────────────────────────────────────────────────────────────────────
\subsection{The genuinely open question: spatial homogenization}
\label{sec:cy_temporal}
%% ────────────────────────────────────────────────────────────────────────────

Because the $m=0$ channel is the leading one, and because the adjoint
shortcut of Section~\ref{sec:cy_adjoint} places no restriction on the
forcing's spatial shape -- a steady (non-oscillating) correction is exactly
as valid an input to it as any other -- the machinery built here is a
reasonable tool to bring to the question that actually remains open (this
specific reduction has not itself been carried through to a checked
assertion, and is flagged as such rather than presented as settled): replacing a
spatially varying $\eta(x,\theta)$ with a single scalar
$\mathrm{Oh}_{\mathrm{eff}}(t)$ is an uncontrolled effective-medium
approximation unless the spatial variation is itself small. Checked
directly, using the period-averaged shear-rate shape from
Section~\ref{sec:cy_strain}--\ref{sec:cy_period_pi} over the drop's volume
at $\mathrm{Oh}_\infty$, $l=2$: the shape varies by a factor of
$\approx 2.25$ between its minimum and maximum over $(x,\theta)$ ($-50\%$ to
$+13\%$ around the mean). Propagated through the Carreau-Yasuda exponent
$a\approx0.743$, this is roughly a factor-of-two spread in local viscosity
across the drop at a single instant -- smaller than the multiple-order-of-
magnitude range $\mathrm{Oh}_{\mathrm{eff}}$ traverses over an impact
\emph{in time}, but not small in the sense a first-order correction needs
to be uncontroversial.

Legendre-projecting the (properly $m=0$, period-averaged, not single-phase)
shear-rate correction shape at a representative radius confirms the
fractional-power leakage mechanism anticipated on general grounds: a
\emph{single} active mode ($l=2$) spontaneously forces every even
spherical-harmonic degree ($l'=0,4,6,8$ all genuinely nonzero, odd $l'$
zero to $10^{-11}$ by a verified $\cos\theta$-parity argument), because
raising a finite Legendre polynomial to a non-integer power does not, in
general, stay within its own degree.

\begin{keybox}[What this section establishes, precisely]
\begin{enumerate}
  \item The classical $O(\epsilon^3)$ Carreau expansion of
        Section~\ref{sec:carreau} does not extend to this fluid's fitted,
        non-integer Carreau-Yasuda exponent: no small-amplitude anchor,
        at rest, at infinite shear, or in between, supports a uniformly
        valid linear correction to $\lambda_l(\mathrm{Oh})$ across the
        modes this solver resolves.
  \item Reid's exact (non-perturbative, continuation-solved)
        $\lambda_l(\mathrm{Oh})$, $\omega_l^2(\mathrm{Oh})$ should be
        evaluated directly at whatever $\mathrm{Oh}_{\mathrm{eff}}(t)$ a
        closure provides -- never linearized -- since this repo already
        implements a fast, tabulated version of exactly that evaluation.
  \item The genuinely open approximation in the existing heuristic scheme
        is the collapse of a spatially varying $\eta(x,\theta,t)$ onto a
        single scalar $\mathrm{Oh}_{\mathrm{eff}}(t)$, quantified here at
        roughly a factor of two in local viscosity across the drop at a
        representative instant -- a real, moderate, but not yet
        controlled error.
  \item Three general, CAS-verified results survive independently of any
        of the above: the adjoint sensitivity shortcut
        (Eq.~\ref{eq:adjoint_shortcut}), the viscous strain-rate tensor
        (Eq.~\ref{eq:strain_tensor}), and the period-$\pi$ lemma. None
        depend on this fluid's specific parameters, and all three are the
        correct starting point for a future, genuinely first-principles
        treatment -- most plausibly a self-consistent, spatially resolved
        solve (using the adjoint machinery as a Newton/Fr\'echet-derivative
        iteration step, not a one-shot small-parameter expansion) rather
        than any further attempt at a closed-form amplitude equation.
\end{enumerate}
\end{keybox}

\begin{notebox}[Scope]
This section does not implement the self-consistent spatially resolved
solve it recommends -- that is future work. It also does not revisit the
$m=2$ parametric channel identified in
Section~\ref{sec:cy_period_pi}, nor extend the strain-tensor/leakage
calculation beyond $l=2$ and a single representative radius. What it does
establish, to the same CAS-verified standard as every other section of this
document -- every numbered claim above has a corresponding assertion in
\texttt{julia/derivations/carreauYasuda\_firstprinciples\_derivation.jl},
including a live cross-check against a running \texttt{solve\_drop!} trace
(the pinned-intermediate-Oh finding of
Section~\ref{sec:cy_anchors}) -- is precisely which of the natural next
steps work, which do not, and why -- for this fluid, not as a general
verdict on Carreau-Yasuda drops.
\end{notebox}


%% ════════════════════════════════════════════════════════════════════════════
\section{Extension to Viscoelastic Drops: the Oldroyd-B Model}
\label{sec:oldroydB}
%% ════════════════════════════════════════════════════════════════════════════

Every result in Sections 1--9 assumes the viscous stress is instantaneously proportional to the rate of strain,
$\boldsymbol{\tau} = 2\mu\,\mathbf{e}$. Real polymer solutions violate this.
A drop of polyethylene oxide (PEO) in water, a silicone oil of moderate
molecular weight, or a biological fluid such as mucus all retain a
\emph{memory} of past deformations. Stress does not vanish the instant
deformation stops; it relaxes on a characteristic timescale set by the
polymer microstructure.

The \textbf{Oldroyd-B model} is the simplest constitutive law that captures
this memory and remains analytically tractable. It is the
canonical model for dilute polymer solutions: a Newtonian solvent (viscosity
$\mu_s$) in which elastic polymer dumbbells are suspended (contributing
additional viscosity $\mu_p$, with $\mu = \mu_s + \mu_p$ the total). The
model has two timescales: a \emph{relaxation time} $\lambda_1$ (how long
polymer stress persists after deformation stops) and a \emph{retardation
time} $\lambda_2 < \lambda_1$ (how long the combined fluid takes to respond
to a new deformation). Setting $\lambda_2 = 0$ reduces Oldroyd-B to the
Maxwell model; setting $\lambda_1 = \lambda_2$ recovers a Newtonian fluid.

Reid's characteristic equation \eqref{eq:char_eq} survives generalisation
to Oldroyd-B. The left-hand side $\alpha^4/q^4 + 1$ is unchanged. The
right-hand side is Reid's right-hand side with $q$ replaced by a modified
wavenumber $q_*$, a rational function of $q^2$. This rational structure
introduces a new type of pole, the \emph{retardation pole}, which captures
additional eigenvalues beyond Reid's fundamental pair and eventually causes
the two-number $\{A_l, D_l\}$ picture to break down.

%% ────────────────────────────────────────────────────────────────────────────
\subsection{The Oldroyd-B constitutive law}
%% ────────────────────────────────────────────────────────────────────────────

\subsubsection{Physical motivation: the mechanical analog}

\begin{notebox}[Spring-dashpot picture of Oldroyd-B]
Imagine the fluid as two \emph{parallel} stress-carrying branches:
\begin{itemize}
  \item \textbf{Solvent branch:} a pure dashpot of viscosity $\mu_s$.
        Stress responds instantaneously to strain rate: $\boldsymbol{\tau}_s
        = 2\mu_s\,\mathbf{e}$. No memory.
  \item \textbf{Polymer branch:} a Hookean spring (elastic modulus
        $G = \mu_p/\lambda_1$) in \emph{series} with a dashpot of
        viscosity $\mu_p$. The spring stores elastic energy when the
        dumbbell is stretched; the dashpot relaxes that energy on timescale
        $\lambda_1 = \mu_p/G$.
\end{itemize}
The total stress is $\boldsymbol{\tau} = \boldsymbol{\tau}_s +
\boldsymbol{\tau}_p$. The retardation time
\begin{equation}
  \lambda_2 = \lambda_1\,\frac{\mu_s}{\mu_s + \mu_p} = \beta_s\,\lambda_1
\end{equation}
is the timescale on which the \emph{combined} system (solvent + polymer)
first responds to a new deformation: the solvent wants to respond
immediately, but the spring in the polymer arm initially resists,
delaying the total response.
\end{notebox}

\subsubsection{The constitutive equation: two equivalent forms}

In terms of the two raw timescales $(\lambda_1, \lambda_2)$:
\begin{equation}
  \boldsymbol{\tau} + \lambda_1\,\overset{\nabla}{\boldsymbol{\tau}}
  = 2\mu\!\left(\mathbf{e} + \lambda_2\,\overset{\nabla}{\mathbf{e}}\right),
  \qquad 0 \leq \lambda_2 < \lambda_1,
  \label{eq:OB_raw}
\end{equation}
where $\overset{\nabla}{(\cdot)}$ is the upper-convected time derivative:
\begin{equation}
  \overset{\nabla}{\mathbf{A}}
  = \frac{\partial\mathbf{A}}{\partial t}
  + (\bu\cdot\nabla)\mathbf{A}
  - (\nabla\bu)\cdot\mathbf{A}
  - \mathbf{A}\cdot(\nabla\bu)^\top.
  \label{eq:ucd}
\end{equation}
In terms of the retardation ratio $\beta_s \equiv \lambda_2/\lambda_1$:
\begin{equation}
  \boldsymbol{\tau} + \lambda_1\,\overset{\nabla}{\boldsymbol{\tau}}
  = 2\mu\!\left(\mathbf{e} + \beta_s\lambda_1\,\overset{\nabla}{\mathbf{e}}\right),
  \qquad \beta_s \in [0,1),
  \label{eq:OB_beta}
\end{equation}
with $\beta_s = \mu_s/(\mu_s+\mu_p)$ the solvent viscosity fraction.
The two forms \eqref{eq:OB_raw} and \eqref{eq:OB_beta} are identical;
we carry both throughout because each makes different limits transparent.

\begin{notebox}[Upper-convected derivative at linearised order]
All fields are $O(\epsilon)$. In the upper-convected derivative
\eqref{eq:ucd}, the terms $(\bu\cdot\nabla)\mathbf{A}$,
$(\nabla\bu)\cdot\mathbf{A}$, and $\mathbf{A}\cdot(\nabla\bu)^\top$ are
each products of two $O(\epsilon)$ quantities, hence $O(\epsilon^2)$.
At linear order:
\[
  \overset{\nabla}{\boldsymbol{\tau}} \;\longrightarrow\;
  \frac{\partial\boldsymbol{\tau}}{\partial t} + O(\epsilon^2), \qquad
  \overset{\nabla}{\mathbf{e}} \;\longrightarrow\;
  \frac{\partial\mathbf{e}}{\partial t} + O(\epsilon^2).
\]
We record the full upper-convected form for completeness and frame-
indifference, but drop the $O(\epsilon^2)$ terms immediately.
\end{notebox}

\subsubsection{Parameter limits}

\begin{center}
\begin{tabular}{lllll}
  \toprule
  Limit & Condition & Model & $\beta_s$ & \\
  \midrule
  Maxwell    & $\lambda_2 = 0$        & Single-relaxation memory & $0$     & \\
  Newtonian  & $\lambda_1 = \lambda_2$& Viscous, viscosity $\mu$ & $1$     & \\
  Solvent    & $\lambda_1\to 0$ at fixed $G=\mu_p/\lambda_1$, $\mu_s$
                                      & Newtonian, viscosity $\mu_s$ & $\to 1$ & \\
  Elastic solid & $\lambda_1\to\infty$, $\beta_s=0$, $G=\mu_p/\lambda_1$ fixed ($\mu_p\to\infty$)
                                      & Hookean solid, modulus $G$ & $0$ & \\
  \bottomrule
\end{tabular}
\end{center}

%% ────────────────────────────────────────────────────────────────────────────
\subsection{Linearisation and effective viscosity}
%% ────────────────────────────────────────────────────────────────────────────

The linearisation strategy in this section follows Zrni\'{c} \& Brenn
\cite{zirnic2024}, whose Eq.\,(3.3) establishes the same effective-viscosity
replacement $\mu \to \mu_{\mathrm{eff}}(\sigma)$ for the free-surface Oldroyd-B
problem.

\begin{derivbox}[Linearised Oldroyd-B in the frequency domain]
With time dependence $e^{-\sigma t}$ (our convention: $\mathrm{Re}(\sigma)>0$
for decay), $\partial_t \to -\sigma$. Dropping $O(\epsilon^2)$ upper-convected
terms:
\[
  (1 - \lambda_1\sigma)\,\boldsymbol{\tau}
  = 2\mu\,(1 - \lambda_2\sigma)\,\mathbf{e}.
\]
Dividing through by $(1 - \lambda_1\sigma)$ (valid generically, since
$\sigma = 1/\lambda_1$ is a special value examined separately):
\begin{equation}
  \boldsymbol{\tau} = 2\mu_{\mathrm{eff}}(\sigma)\,\mathbf{e},
  \qquad
  \mu_{\mathrm{eff}}(\sigma) = \mu\,\frac{1-\lambda_2\sigma}{1-\lambda_1\sigma}.
  \label{eq:mueff_OB}
\end{equation}
In both parameterisations simultaneously:
\begin{equation}
  \mu_{\mathrm{eff}}(\sigma)
  = \mu\,\frac{1-\lambda_2\sigma}{1-\lambda_1\sigma}
  = \mu\,\frac{1-\beta_s\lambda_1\sigma}{1-\lambda_1\sigma}.
  \label{eq:mueff_both}
\end{equation}
The effective kinematic viscosity is $\nu_{\mathrm{eff}} =
\mu_{\mathrm{eff}}/\rho$.
\end{derivbox}

\begin{notebox}[Analytic structure of $\mu_{\mathrm{eff}}(\sigma)$]
As a function of $\sigma$, $\mu_{\mathrm{eff}}$ is a ratio of two linear
factors:
\begin{itemize}
  \item \textbf{Zero} at $\sigma = 1/\lambda_2$: solvent and polymer stresses
        cancel in the tangential direction. ($\lambda_2 = 0$ pushes this
        zero to infinity (the Maxwell model has no zero.)
  \item \textbf{Pole} at $\sigma = 1/\lambda_1$: polymer stress cannot relax;
        effective viscosity diverges (solid-like response).
\end{itemize}
Setting $\lambda_2 = 0$: $\mu_{\mathrm{eff}} = \mu/(1-\lambda_1\sigma)$,
recovering the Maxwell effective viscosity. Setting $\lambda_1 = \lambda_2$:
$\mu_{\mathrm{eff}} = \mu$, recovering Newtonian. The Oldroyd-B model
interpolates between these two limits.
\end{notebox}

\subsubsection{The two Deborah numbers}

Referenced to the inviscid oscillation frequency $\sigma_{l;0}$:
\begin{equation}
  De_1 \equiv \lambda_1\,\sigma_{l;0},
  \qquad
  De_2 \equiv \lambda_2\,\sigma_{l;0} = \beta_s\,De_1.
  \label{eq:De_defs}
\end{equation}
$De_j$ is the ratio of the fluid's $j$th memory timescale to the drop's
natural oscillation period.

\begin{notebox}[Convention note]
Zrni\'{c} \& Brenn \cite{zirnic2024} define their Deborah numbers referenced
to the capillary timescale $(T_1/(\rho R^3))^{1/2}$, whereas here they are
referenced to $\sigma_{l;0} = [l(l-1)(l+2)]^{1/2}(T_1/(\rho R^3))^{1/2}$.
Consequently $De_j^{\text{here}} = \sqrt{l(l-1)(l+2)}\,De_j^{\text{Z\&B}}$;
numerical values are mode-dependent and should not be compared directly.
\end{notebox} When $De_1 \ll 1$, the polymer relaxes many times
per cycle and the drop behaves as Newtonian. When $De_1 \sim 1$, elasticity
and surface tension compete on equal footing.

\begin{derivbox}[Expressing $\lambda_j\sigma$ in terms of $q^2$ and $\alpha^2$]
Recall $\sigma = q^2\nu/R^2$ and $\sigma_{l;0} = \alpha^2\nu/R^2$.
Therefore:
\begin{equation}
  \lambda_j\sigma
  = \lambda_j\cdot\frac{q^2\nu}{R^2}
  = \frac{\lambda_j\sigma_{l;0}}{\alpha^2\nu/R^2}\cdot\frac{q^2\nu}{R^2}
  = \frac{De_j\cdot q^2}{\alpha^2}.
  \label{eq:lamsig}
\end{equation}
Substituting into \eqref{eq:mueff_OB}:
\begin{equation}
  \mu_{\mathrm{eff}} = \mu\,\frac{\alpha^2 - De_2\,q^2}{\alpha^2 - De_1\,q^2},
  \qquad
  \nu_{\mathrm{eff}} = \nu\,\frac{\alpha^2 - De_2\,q^2}{\alpha^2 - De_1\,q^2}.
  \label{eq:nueff_dimless}
\end{equation}
All dimensional quantities have cancelled. The effective viscosity is a
function of $q^2/\alpha^2$ only.
\end{derivbox}

%% ────────────────────────────────────────────────────────────────────────────
\subsection{Boundary condition audit}
%% ────────────────────────────────────────────────────────────────────────────

We have a free-surface drop in air. We examine each of the three boundary
conditions.

\subsubsection{BC1 (kinematic): unchanged}

The kinematic condition equates the fluid radial velocity at $x=1$ to the
rate of change of the surface position. It is a purely geometric statement
involving only the \emph{velocity field}, not the stress. The Oldroyd-B
constitutive law does not alter the velocity field structure. Therefore:
\begin{equation}
  U(1) = -1. \tag{BC1, unchanged}
\end{equation}

\subsubsection{BC2 (tangential stress): unchanged in form}

The free-surface condition $\tau_{r\theta}|_{x=1} = 0$ becomes:
\[
  \mu_{\mathrm{eff}}(\sigma)\!\left[\,r\,\frac{\partial}{\partial r}
  \!\left(\frac{u_\theta}{r}\right)
  + \frac{1}{r}\frac{\partial u_r}{\partial\theta}\,\right]_{r=R} = 0.
\]
Since $\mu_{\mathrm{eff}} \neq 0$ generically, we divide through. The
condition reduces to exactly the same operator equation on $U$:
\begin{equation}
  \left[\,\frac{\dd^2}{\dd x^2}
  - \frac{2}{x}\frac{\dd}{\dd x}
  + \frac{l(l+1)}{x^2}\,\right]U\,\bigg|_{x=1} = 0. \tag{BC2, unchanged}
\end{equation}

\begin{notebox}[The zero of $\mu_{\mathrm{eff}}$ at $\sigma = 1/\lambda_2$]
At $\sigma = 1/\lambda_2$, the effective viscosity $\mu_{\mathrm{eff}} \to 0$
and BC2 becomes $0 = 0$ for any $U$: it imposes no constraint.
The frequency $\sigma = 1/\lambda_2$ is the \emph{retardation pole} of
$q_*^2$: as $w \to w_{\mathrm{ret}} = \alpha^2/De_2$, $q_*^2 \to \infty$.
It is not itself a root of the characteristic equation $F(w) = 0$. The actual retardation roots instead lie in a neighbourhood of $w_{\mathrm{ret}}$ (see Section~\ref{sec:OB_roots}).
At those roots BC2 must be handled by a limiting argument. In practice these
roots decay at rate $1/\lambda_2$, which is faster than the capillary
timescale when $De_2 < 1$, and they are discarded as overtones. For
$De_2 > 1$ a more careful treatment is required.
\end{notebox}

\subsubsection{BC3 (normal stress): viscoelastic correction check}

\begin{derivbox}[First normal stress difference at $O(\epsilon)$ for Oldroyd-B]
For the Oldroyd-B family, the first normal stress difference is:
\[
  N_1 = \tau_{rr} - \tau_{\theta\theta}.
\]
At the linearised level, $\boldsymbol{\tau} = 2\mu_{\mathrm{eff}}\,\mathbf{e}$
with $\mathbf{e} = O(\epsilon)$, so $N_1 = O(\epsilon)$. Does this $O(\epsilon)$
quantity add a new term to BC3?

The normal stress balance at a free surface equates the jump in
\emph{radial} normal stress to the Laplace pressure:
\[
  -p_{rr}\big|_{\text{inside}} = T_1\!\left(\frac{1}{R_1}+\frac{1}{R_2}\right).
\]
The radial normal stress inside is $-p_{rr} = p + \delta p - \tau_{rr}$,
with $\tau_{rr} = 2\mu_{\mathrm{eff}}\,\partial u_r/\partial r$. This is
precisely the viscous term already present in BC3 \eqref{eq:BC3}, now with
$\mu \to \mu_{\mathrm{eff}}$. The remaining part $N_1 = \tau_{rr} -
\tau_{\theta\theta}$ measures the \emph{anisotropy} of the stress tensor,
but it does not enter BC3 at $O(\epsilon)$.
The reason is that the base state is a quiescent sphere with
$\boldsymbol{\tau}_0 = 0$, so the $O(\epsilon)$ tilt of the outward normal
from $\hat{\mathbf{r}}$ introduces only an $O(\epsilon)$ correction to the
normal direction times an $O(\epsilon^0) = 0$ base tangential traction,
contributing nothing at $O(\epsilon)$.
Therefore $N_1$ contributes nothing extra to BC3 at $O(\epsilon)$.

At $O(\epsilon^2)$, the surface is deformed and the mismatch between normal
and tangential directions at the perturbed surface generates an
$O(\epsilon^2)$ correction from $N_1$ to the effective surface tension.
This is an open problem noted in Section~\ref{sec:OB_open}.
\end{derivbox}

The normal stress boundary condition is:
\begin{equation}
  (l-1)(l+2)\frac{T_1}{\rho R}
  = P_0 - 2\nu_{\mathrm{eff}}(\sigma)\,\sigma\!\left[\,U'(1) - 2U(1)\,\right].
  \tag{BC3$'$}
  \label{eq:BC3_OB}
\end{equation}
This is identical to the Newtonian BC3 \eqref{eq:BC3} with
$\nu \to \nu_{\mathrm{eff}}$ from \eqref{eq:nueff_dimless}.

%% ────────────────────────────────────────────────────────────────────────────
\subsection{Modified governing ODE}
%% ────────────────────────────────────────────────────────────────────────────

\subsubsection{The modified wavenumber $q_*$}

\begin{derivbox}[Definition and derivation of $q_*^2$]
In Reid's Newtonian derivation, the parameter $q^2 = \sigma R^2/\nu$ emerged
when the $r$-momentum equation was multiplied through by $R/(\nu\sigma)$.
For Oldroyd-B, the same step with $\nu \to \nu_{\mathrm{eff}}$ gives:
\begin{equation}
  q_*^2 \;\equiv\; \frac{\sigma R^2}{\nu_{\mathrm{eff}}(\sigma)}
  = q^2\,\frac{1-\lambda_1\sigma}{1-\lambda_2\sigma}.
  \label{eq:qstar_raw}
\end{equation}
Using \eqref{eq:lamsig} to express $\lambda_j\sigma$ in terms of $q^2$
and $\alpha^2$:
\begin{equation}
  \boxed{
    q_*^2 = q^2\cdot\frac{\alpha^2 - De_1\,q^2}{\alpha^2 - De_2\,q^2}.
  }
  \label{eq:qstar_OB}
\end{equation}
In both parameterisations:
\begin{equation}
  q_*^2 = q^2\cdot\frac{\alpha^2 - De_1\,q^2}{\alpha^2 - \beta_s De_1\,q^2}.
  \label{eq:qstar_both}
\end{equation}
\end{derivbox}

\begin{notebox}[Structural difference from the Maxwell model]
For the Maxwell fluid, $q_*^2 = q^2(1 - De_1 q^2/\alpha^2) =
q^2(\alpha^2 - De_1 q^2)/\alpha^2$: the denominator is the constant
$\alpha^2$. For Oldroyd-B, the denominator is $\alpha^2 - De_2 q^2$,
which \emph{vanishes} at $q^2 = \alpha^2/De_2$, i.e.\ at
$\sigma = \sigma_{l;0}/De_2 = 1/\lambda_2$. As $q^2 \to \alpha^2/De_2$:
$q_*^2 \to \infty$. This is the \textbf{retardation pole}: a singularity
in the effective wavenumber absent from Maxwell or Newtonian problems. All qualitative differences between Oldroyd-B and Maxwell originate in this pole.
\end{notebox}

\begin{notebox}[Limits of $q_*^2$]
\begin{itemize}
  \item $De_2 = 0$ (Maxwell): $q_*^2 = q^2(\alpha^2 - De_1 q^2)/\alpha^2$.
        Retardation pole disappears.
  \item $De_1 = De_2$ (Newtonian): numerator equals denominator,
        $q_*^2 = q^2$.
  \item $q^2 = \alpha^2/De_1$ (relaxation zero): numerator $= 0$,
        $q_*^2 = 0$. At $\sigma = 1/\lambda_1$ the polymer is fully
        relaxed; it contributes zero stress, as if absent.
  \item $q^2 = \alpha^2/De_2$ (retardation pole): denominator $= 0$,
        $q_*^2\to\infty$. At $\sigma = 1/\lambda_2$ the effective
        viscosity vanishes and the Bessel argument diverges.
\end{itemize}
Since $De_2 < De_1$, the retardation pole $\alpha^2/De_2 > \alpha^2/De_1$
lies to the right of the relaxation zero on the positive real $q^2$-axis.
\end{notebox}

\subsubsection{The ODE for $U(x)$ under Oldroyd-B}

\begin{derivbox}[The pressure field is unchanged]
The pressure perturbation satisfies $\nabla^2(\delta p/\rho) = 0$ regardless
of $\nu_{\mathrm{eff}}$. To see this: take the divergence of the linearised
momentum equation
$-\sigma\bu = -\nabla(\delta p/\rho) + \nu_{\mathrm{eff}}\nabla^2\bu$
and use incompressibility $\nabla\cdot\bu = 0$. Since $\nu_{\mathrm{eff}}$
is a scalar constant (it depends on $\sigma$ but not on $\mathbf{x}$), we get
$\nabla^2(\delta p/\rho) = \nu_{\mathrm{eff}}\nabla^2(\nabla\cdot\bu) = 0$.
The solution $\delta p/\rho = \epsilon P_0 x^l Y_l^m$ is unchanged.
\end{derivbox}

Following Section~5 with the substitution $\nu \to \nu_{\mathrm{eff}}$
and hence $q^2 \to q_*^2$:

\begin{keybox}[ODE for $U(x)$ under Oldroyd-B]
\begin{equation}
  \left[\,\frac{\dd^2}{\dd x^2} - \frac{l(l+1)}{x^2} + q_*^2\,\right]U(x)
  = q_*^2\,\Pi_0\,x^{l+1},
  \label{eq:ODE_OB}
\end{equation}
with $q_*^2$ given by \eqref{eq:qstar_OB}. The general solution is:
\begin{equation}
  U(x) = C\,x\,j_l(q_* x) + \Pi_0\,x^{l+1}.
  \label{eq:U_OB}
\end{equation}
This is Reid's ODE and general solution \eqref{eq:ODE_U}--\eqref{eq:U_general}
with the single replacement $q \to q_*$.
\end{keybox}

%% ────────────────────────────────────────────────────────────────────────────
\subsection{The Oldroyd-B characteristic equation}
%% ────────────────────────────────────────────────────────────────────────────

\subsubsection{The ratio cancellation}

Before deriving the characteristic equation, we establish a key algebraic
fact that determines its structure.

\begin{derivbox}[The ratio $\alpha_*^2/q_*^2$ is Newtonian]
Define, in analogy with \eqref{eq:alpha_def}:
\begin{equation}
  \alpha_*^2 \;\equiv\; \frac{\sigma_{l;0}R^2}{\nu_{\mathrm{eff}}}
  = \alpha^2\,\frac{1-\lambda_1\sigma}{1-\lambda_2\sigma}
  = \alpha^2\cdot\frac{\alpha^2-De_1\,q^2}{\alpha^2-De_2\,q^2}.
  \label{eq:alphastar_OB}
\end{equation}
Computing the ratio $\alpha_*^2/q_*^2$:
\[
  \frac{\alpha_*^2}{q_*^2}
  = \frac{\displaystyle\alpha^2\cdot\frac{\alpha^2-De_1 q^2}{\alpha^2-De_2 q^2}}
         {\displaystyle q^2\cdot\frac{\alpha^2-De_1 q^2}{\alpha^2-De_2 q^2}}
  = \frac{\alpha^2}{q^2}.
\]
The Oldroyd-B factors cancel exactly in the ratio. Therefore:
\begin{equation}
  \frac{\alpha_*^4}{q_*^4} = \frac{\alpha^4}{q^4}.
  \label{eq:ratio_cancel}
\end{equation}
\end{derivbox}

\begin{notebox}[Why the cancellation occurs]
The ratio $\alpha^4/q^4 = \sigma_{l;0}^2/\sigma^2$ is a pure ratio of
frequencies. It is independent of viscosity by construction, since both
$\alpha^2 \propto 1/\nu$ and $q^2 \propto 1/\nu$ rescale identically when
$\nu \to \nu_{\mathrm{eff}}$. The cancellation \eqref{eq:ratio_cancel} holds
for \emph{any} constitutive model that linearises to $\boldsymbol{\tau} =
2\mu_{\mathrm{eff}}(\sigma)\,\mathbf{e}$ with a scalar $\mu_{\mathrm{eff}}$.
It is a universal feature of the linearised problem, not specific to
Oldroyd-B.
\end{notebox}

\subsubsection{Eliminating $T_1$}

\begin{derivbox}[Modified rescaling]
Multiplying BC3$'$ \eqref{eq:BC3_OB} through by $lR^2/\nu_{\mathrm{eff}}^2$
and using $(l-1)(l+2)T_1/(\rho R) = \sigma_{l;0}^2 R^2/l$:
\[
  \underbrace{\frac{\sigma_{l;0}^2 R^4}{\nu_{\mathrm{eff}}^2}}_{\alpha_*^4}
  = \underbrace{\frac{\sigma^2 R^4}{\nu_{\mathrm{eff}}^2}}_{q_*^4}\Pi_0
  - 2l\underbrace{\frac{\sigma R^2}{\nu_{\mathrm{eff}}}}_{q_*^2}\!\left[U'(1)-2U(1)\right].
\]
Hence:
\begin{equation}
  \alpha_*^4 = q_*^2\!\left\{\,q_*^2\Pi_0
  - 2l\!\left[U'(1) - 2U(1)\right]\right\}.
  \label{eq:alpha4star_OB}
\end{equation}
$T_1$ and $\rho$ have cancelled completely. By \eqref{eq:ratio_cancel},
$\alpha_*^4 = q_*^4\cdot\alpha^4/q^4$, so the left-hand side of the
characteristic equation will be $\alpha^4/q^4$, identical to the Newtonian
case.
\end{derivbox}

\subsubsection{Applying BC1 and BC2 and assembling the characteristic equation}

The algebra of Section~7.2 carries through verbatim with $q \to q_*$:

\paragraph{BC1:} $U(1) = C\,j_l(q_*) + \Pi_0 = -1.$

\paragraph{BC2:} Evaluating \eqref{eq:BC2} on \eqref{eq:U_OB}:
\[
  C\!\left[(2l-q_*^2)\,j_l(q_*) - 2q_*\,j_l'(q_*) + 2(l^2-1)\,j_l(q_*)\right]
  + 2(l^2-1)\Pi_0 = 0.
\]

\paragraph{Solving for $C$ and $\Pi_0$:}
Substituting $\Pi_0 = -1 - C j_l(q_*)$:
\begin{equation}
  C = \frac{2(l-1)(l+1)}{(2l-q_*^2)\,j_l(q_*) - 2q_*\,j_l'(q_*)}.
  \label{eq:C_OB}
\end{equation}
Using the Bessel derivative identity $q_* j_l'(q_*)/j_l(q_*) =
l - q_* Q_{l+1/2}(q_*)$ and the Bessel recurrence to simplify, then
substituting $C$, $\Pi_0$, $U'(1)$, and $U(1)$ into \eqref{eq:alpha4star_OB}:

\begin{keybox}[Oldroyd-B Characteristic Equation]
\begin{equation}
  \frac{\alpha^4}{q^4} + 1
  = \frac{2(l-1)}{q_*^2}
    \left[\,l + (l+1)\,\frac{q_* - 2l\,Q_{l+1/2}(q_*)}{q_* - 2\,Q_{l+1/2}(q_*)}\,\right],
  \label{eq:char_OB}
\end{equation}
where
\begin{equation}
  q_*^2 = q^2\cdot\frac{\alpha^2-De_1\,q^2}{\alpha^2-De_2\,q^2}
  = q^2\cdot\frac{\alpha^2-De_1\,q^2}{\alpha^2-\beta_s De_1\,q^2},
  \label{eq:qstar_OB_final}
\end{equation}
$Q_{l+1/2}(q_*) = J_{l+3/2}(q_*)/J_{l+1/2}(q_*)$, and the parameters
$\alpha^2 = \sigma_{l;0}R^2/\nu$, $De_1 = \lambda_1\sigma_{l;0}$,
$De_2 = \lambda_2\sigma_{l;0} = \beta_s De_1$ are real constants.

\smallskip\noindent Limits:
$De_1 = 0$ (or $\beta_s = 1$): $q_* = q$, recovers Reid \eqref{eq:char_eq}.\\
$De_2 = 0$: recovers the Maxwell characteristic equation.\\
$De_1 = De_2 \neq 0$: $q_* = q$, again recovers Reid \eqref{eq:char_eq}.\\[2pt]
This equation coincides with Zrni\'{c} \& Brenn \cite{zirnic2024} Eqs.\,(3.11)--(3.12).
\end{keybox}

\begin{notebox}[What changed from Reid's equation]
\begin{enumerate}
  \item The left-hand side $\alpha^4/q^4 + 1$ is \emph{identical} to Reid's.
        (Consequence of the ratio cancellation \eqref{eq:ratio_cancel}.)
  \item The right-hand side is Reid's right-hand side with $q \to q_*$,
        where $q_*$ is the rational function \eqref{eq:qstar_OB_final}.
  \item The entire Oldroyd-B physics is encoded in $De_1$ and $\beta_s$
        (two extra real parameters beyond Reid's $\alpha^2$).
  \item The rational structure of $q_*^2$ introduces a new pole
        at $q^2 = \alpha^2/De_2$ with no Newtonian or Maxwell analogue.
\end{enumerate}
\end{notebox}

%% ────────────────────────────────────────────────────────────────────────────
\subsection{Structure of solutions}
\label{sec:OB_roots}
%% ────────────────────────────────────────────────────────────────────────────

\subsubsection{The characteristic equation as a condition in $q^2$}

Let $w = q^2$. Define the rational map:
\begin{equation}
  R(w) \;\equiv\; q_*^2
  = w\cdot\frac{\alpha^2 - De_1 w}{\alpha^2 - De_2 w}.
  \label{eq:R_map}
\end{equation}
Then \eqref{eq:char_OB} is a transcendental equation $F(w) = 0$ in the single
complex variable $w$, with $\alpha^2$, $De_1$, $De_2$ as real parameters.
Its roots $w_j$ give eigenvalues $\sigma_j = w_j\nu/R^2$.

\subsubsection{Two types of poles in $F(w)$}
\label{sec:two_pole_types}

The function $F(w)$ has poles wherever the right-hand side of
\eqref{eq:char_OB} diverges. There are two distinct sources:

\paragraph{Type 1: Bessel poles.}
The right-hand side of \eqref{eq:char_OB} has poles wherever
$J_{l+1/2}(q_*) = 0$, i.e.\ at the zeros $z_1^* < z_2^* < \cdots$ of the
Bessel function in the $q_*$-plane. These satisfy $q_* = z_n^*$, i.e.\
$R(w) = (z_n^*)^2$: for each $n$, this is a quadratic (in $w$) algebraic
equation with up to two solutions in the $w$-plane. For small $De_1$ the rational map $R(w)$ is a small perturbation of the
Newtonian case, and a heuristic application of the argument principle suggests
that $F(w) = 0$ retains exactly two roots per strip, as in the Newtonian
problem; a rigorous contour argument for general $De_1$ is open (see
\S\ref{sec:OB_open}). The first strip ($n=1$) captures the
\textbf{dominant pair}, governing the fundamental capillary oscillation.

\paragraph{Type 2: Retardation pole.}
The factor $\alpha^2 - De_2 w$ in the denominator of $R(w)$ vanishes at:
\begin{equation}
  w_{\mathrm{ret}} = \frac{\alpha^2}{De_2},
  \qquad
  \sigma_{\mathrm{ret}} = \frac{\sigma_{l;0}}{De_2} = \frac{1}{\lambda_2}.
  \label{eq:ret_pole}
\end{equation}
As $w \to w_{\mathrm{ret}}$, $q_* \to \infty$ and $Q_{l+1/2}(q_*)$ oscillates
rapidly. Near this singularity the characteristic equation
\eqref{eq:char_OB} can be satisfied, capturing one or two
\textbf{retardation roots}. These correspond to modes decaying at rate
$\sim 1/\lambda_2$. Since $\lambda_2 < \lambda_1$, the retardation mode
always decays faster than the Maxwell-type relaxation mode; whether it is
faster or slower than the fundamental capillary pair depends on $De_2$.

\begin{notebox}[Physical meaning of the retardation roots]
At $\sigma = 1/\lambda_2$, the solvent and polymer contributions to the
tangential stress cancel: $\mu_{\mathrm{eff}}(1/\lambda_2) = 0$. The
retardation roots describe a mode in which the polymer network, instead of
resisting the flow, briefly \emph{drives} it, before the solvent dissipates
the energy. This phenomenon has no Newtonian analogue; it requires two fluid components. Physically, it corresponds to the elastic recoil
of the polymer network overshooting the equilibrium and being damped by the
solvent on the retardation timescale $\lambda_2$.
\end{notebox}

\subsubsection{Dominance condition for the fundamental pair}

The fundamental pair of roots governs observable drop oscillation provided
the retardation roots are faster-decaying (subdominant):
\begin{equation}
  \mathrm{Re}(\sigma_{1,2}) < \frac{1}{\lambda_2}
  \quad\Longleftrightarrow\quad
  De_2\cdot\mathrm{Re}(\beta_{1,2}) < 1,
  \label{eq:dominance}
\end{equation}
where $\beta_{j} = q_j^2/\alpha^2 = \sigma_j/\sigma_{l;0}$.

When \eqref{eq:dominance} holds, the retardation roots are higher overtones:
they decay faster than the capillary oscillation period and are irrelevant to
the Mola\v{c}ek \& Bush quasi-static model. When \eqref{eq:dominance} fails,
the retardation roots must be retained.

%% ────────────────────────────────────────────────────────────────────────────
\subsection{$A_l$ and $D_l$ for the Oldroyd-B fluid}
%% ────────────────────────────────────────────────────────────────────────────

\subsubsection{Vieta extraction: the two-number picture survives}

Provided \eqref{eq:dominance} holds, the dynamics of mode $l$ is governed by
the two roots $\beta_{1,2}$ of the characteristic equation \eqref{eq:char_OB}
in the first Bessel strip. The Vieta extraction of Section~9 applies without
modification:
\begin{equation}
  A_l^{\mathrm{OB}} = \frac{1}{\beta_1\beta_2},
  \qquad
  D_l^{\mathrm{OB}} = \frac{\beta_1 + \beta_2}{2a\,\beta_1\beta_2},
  \qquad
  a = \mathrm{Oh}\sqrt{l(l-1)(l+2)}.
  \label{eq:AD_OB}
\end{equation}
The difference from the Newtonian case is that $\beta_{1,2}$ are now roots
of \eqref{eq:char_OB} rather than \eqref{eq:char_eq}: they depend on
$(De_1, \beta_s)$ in addition to $(l, \mathrm{Oh})$.

\begin{keybox}[Parameter dependence of $A_l^{\mathrm{OB}}$ and $D_l^{\mathrm{OB}}$]
\begin{equation}
  A_l^{\mathrm{OB}} = A_l^{\mathrm{OB}}(l,\,\mathrm{Oh},\,De_1,\,\beta_s),
  \qquad
  D_l^{\mathrm{OB}} = D_l^{\mathrm{OB}}(l,\,\mathrm{Oh},\,De_1,\,\beta_s).
\end{equation}
Limits:
\begin{itemize}
  \item $De_1 \to 0$ (or $\beta_s \to 1$): recover Mola\v{c}ek \& Bush
        Newtonian values.
  \item $\beta_s \to 0$ at fixed $De_1$: recover Maxwell values (a special
        case of the Oldroyd-B framework).
  \item $De_1 \to \infty$ at fixed $\beta_s$: strongly elastic limit; both
        coefficients must be determined numerically.
\end{itemize}
\end{keybox}

\begin{notebox}[Monotone inequalities between models (conjectured)]
For the same total viscosity $\mu = \mu_s + \mu_p$ and the same Oh,
numerical evidence suggests:
\[
  A_l^{\mathrm{N}} \;\leq\; A_l^{\mathrm{OB}} \;\leq\; A_l^{\mathrm{M}},
  \qquad
  D_l^{\mathrm{M}} \;\leq\; D_l^{\mathrm{OB}} \;\leq\; D_l^{\mathrm{N}}.
\]
Elasticity enhances effective inertia (increases $A_l$) and reduces
instantaneous dissipation (decreases $D_l$). The solvent, present in
Oldroyd-B but absent in Maxwell, partially restores both toward their
Newtonian values. The interpolation appears monotone in $\beta_s$: as
$\beta_s$ increases from 0 (Maxwell) to 1 (Newtonian), $A_l$ decreases
and $D_l$ increases. These inequalities are consistent with physical
intuition but have not been established analytically from the
transcendental characteristic equation; a formal proof is open.
\end{notebox}

\subsubsection{When the two-number picture breaks}

When condition \eqref{eq:dominance} fails, the retardation roots have decay
rates comparable to the fundamental pair. There are then three (or more) roots
of comparable magnitude, and the Mola\v{c}ek \& Bush quadratic ODE for $b_l$
is no longer adequate. In capillary time $\tau = \sigma_{l;0}\,t$ (all quantities
dimensionless), the full description is an integro-differential equation,
\begin{equation}
  A_l^{\mathrm{OB}}\,\ddot{b}_l
  + 2a
    \int_0^\tau K^{(\mathrm{cap})}(\tau-\tau')\,\dot{b}_l(\tau')\,\dd\tau'
  + b_l = 0, \qquad a \equiv \mathrm{Oh}\sqrt{l(l-1)(l+2)},
  \label{eq:integrodiff_OB}
\end{equation}
where dots denote $\dd/\dd\tau$ and the dimensionless OB memory kernel in
capillary time is
\begin{equation}
  K^{(\mathrm{cap})}(\tau) = \beta_s\,\delta(\tau)
  + \frac{1-\beta_s}{De_1}\,e^{-\tau/De_1}.
  \label{eq:OB_kernel_cap}
\end{equation}
Its positive-exponent Laplace transform is
$\tilde{K}^{(\mathrm{cap})}(\beta)
  = \int_0^\infty K^{(\mathrm{cap})}(\tau)\,e^{\beta\tau}\,\dd\tau
  = \nu_{\mathrm{eff}}(\beta\,\sigma_{l;0})/\nu
  = (1 - De_2\,\beta)/(1 - De_1\,\beta)$,
and its zeroth moment is $\int_0^\infty K^{(\mathrm{cap})}\,\dd\tau = \beta_s + (1-\beta_s) = 1$.
In the Newtonian limit ($De_1\to 0$, or equivalently $\beta_s\to 1$):
$K^{(\mathrm{cap})}\to\delta(\tau)$, and the equation reduces to
$A_l^{\mathrm{OB}}\ddot{b}+2a\dot{b}+b=0$.
This is a structured approximation. The convolution kernel replaces the
full Newtonian damping function $D_l(\mathrm{Oh})$ by the zeroth moment
$\tilde{K}^{(\mathrm{cap})}(0)=1$. It captures the frequency-dependent viscosity
$\nu_{\mathrm{eff}}(\beta\sigma_{l;0})/\nu$ but not the Oh-dependence in
Reid's Bessel-function structure. The exact $D_l(\mathrm{Oh})$ is obtained via
Vieta extraction from \eqref{eq:char_OB} at $De_1=0$.
The dimensional kernel $K_{\mathrm{OB}}(s)$ satisfies
$\tilde{K}_{\mathrm{OB}}(\sigma)
  = \int_0^\infty K_{\mathrm{OB}}(s)\,e^{\sigma s}\,\dd s
  = \nu_{\mathrm{eff}}(\sigma)/\nu$:
\begin{equation}
  K_{\mathrm{OB}}(s) = \beta_s\,\delta(s)
  + \frac{1-\beta_s}{\lambda_1}\,e^{-s/\lambda_1},
  \label{eq:OB_kernel}
\end{equation}
related to the capillary kernel by
$K^{(\mathrm{cap})}(\tau) = \sigma_{l;0}^{-1}\,K_{\mathrm{OB}}(\tau/\sigma_{l;0})$.
The $\delta$ term encodes the instantaneous \textbf{solvent} response with
weight $\beta_s = \mu_s/\mu$; the exponential term encodes the polymer memory
with weight $(1-\beta_s) = \mu_p/\mu$ and relaxation time $\lambda_1$.
The two-number approximation corresponds to replacing $K^{(\mathrm{cap})}$ by
its zeroth moment $= 1$ (treating memory as instantaneous), valid when
$De_2 \cdot \mathrm{Re}(\beta_{1,2}) \ll 1$.

%% ────────────────────────────────────────────────────────────────────────────
\subsection{Four clean limits}
%% ────────────────────────────────────────────────────────────────────────────

\begin{center}
\begin{tabular}{lllll}
  \toprule
  Limit & Condition & $q_*^2$ & Char.\ eq.\ & $A_l,D_l$ \\
  \midrule
  Newtonian
    & $De_1 = 0$
    & $q^2$
    & Reid \eqref{eq:char_eq}
    & Mola\v{c}ek \& Bush \\[4pt]
  Maxwell
    & $\beta_s = 0$
    & $q^2\!\left(1 - \tfrac{De_1\,q^2}{\alpha^2}\right)$
    & Khismatullin \& Nadim
    & \eqref{eq:AD_OB}, $\beta_s=0$ \\[4pt]
  Weak elasticity
    & $De_1\ll 1$
    & $q^2[1-(De_1\!-\!De_2)q^2/\alpha^2]$
    & Perturbation of \eqref{eq:char_eq}
    & Analytic correction \\[4pt]
  No retardation
    & $De_2\ll De_1$
    & $\approx q^2(1-De_1 q^2/\alpha^2)$
    & $\approx$ Maxwell
    & $\approx$ Maxwell values \\
  \bottomrule
\end{tabular}
\end{center}

%% ────────────────────────────────────────────────────────────────────────────
\subsection{Open problems and references}
\label{sec:OB_open}
%% ────────────────────────────────────────────────────────────────────────────

\paragraph{What is established.}
\begin{enumerate}
  \item The characteristic equation \eqref{eq:char_OB}--\eqref{eq:qstar_OB_final}
        is exact at $O(\epsilon)$ for a free-surface Oldroyd-B drop in air.
  \item The ratio cancellation \eqref{eq:ratio_cancel} holds for any
        linearisable constitutive model, making the left-hand side of the
        characteristic equation universal.
  \item The dominance condition \eqref{eq:dominance} gives an explicit
        criterion for when the two-number $\{A_l,D_l\}$ extraction is valid.
  \item The integro-differential ODE \eqref{eq:integrodiff_OB}--\eqref{eq:OB_kernel}
        is the natural viscoelastic extension of the Mola\v{c}ek \& Bush quadratic,
        capturing the frequency-dependent damping $\nu_{\mathrm{eff}}/\nu$ when
        the two-number picture breaks down.
\end{enumerate}

\paragraph{What is open.}
\begin{enumerate}
  \item \textbf{Numerical tabulation.} A complete table of
        $A_l^{\mathrm{OB}}(l,\mathrm{Oh},De_1,\beta_s)$ and
        $D_l^{\mathrm{OB}}(l,\mathrm{Oh},De_1,\beta_s)$ analogous to
        Mola\v{c}ek \& Bush Table~1 does not exist in the literature.
        The Newton's method strategy of Section~8 applies directly: for each
        $(De_1,\beta_s)$, initialise from the Newtonian roots and track them
        as $(De_1,\beta_s)$ are increased.
  \item \textbf{Large $De_1$ root structure.} For $De_1 \gg 1$, the numerator
        $\alpha^2 - De_1 q^2$ changes sign for real $q^2 > \alpha^2/De_1$,
        making $q_*^2$ negative there. The physical consequence is that
        Bessel functions $j_l(q_* x)$ become modified spherical Bessel
        functions $i_l(|q_*|x)$, changing the radial profile from
        oscillatory to monotone. A complete argument-principle analysis for
        large $De_1$ does not exist.
  \item \textbf{Second-order normal stress correction.} The $O(\epsilon^2)$
        contribution from $N_1 = \tau_{rr} - \tau_{\theta\theta}$ modifies
        the effective surface tension as
        $T_1 \to T_1^{\mathrm{eff}} = T_1\!\left[1 + \tilde{c}\,\epsilon\,\mathrm{Re}(N_1^{(2)}R/T_1)\right]$
        with $\tilde{c}$ a dimensionless $O(1)$ coefficient (the precise value
        requires the second-order calculation).
        For $\epsilon = O(0.1\text{--}0.3)$ (typical bouncing-drop
        experiments), this correction could reach several percent in $\sigma_{l;0}$.
  \item \textbf{Experimental validation.} No
        experiment has directly measured $A_l$ and $D_l$ for a viscoelastic
        drop and compared to the characteristic equation \eqref{eq:char_OB}.
        PEO-water drops in the range $De_1 \in [0.1, 2]$ and $\beta_s \in
        [0.5, 0.9]$ would be the natural first target.
\end{enumerate}

\paragraph{Key references.}
\begin{itemize}
  \item \textbf{Oldroyd (1950)} \cite{oldroyd1950}: original derivation of
        \eqref{eq:OB_raw}; clearest source for the physical meaning of
        $\lambda_1$, $\lambda_2$, and $\beta_s$.
  \item \textbf{Bird, Armstrong \& Hassager (1987)} \cite{bird1987}:
        dumbbell-model derivation of Oldroyd-B; physical meaning of
        $\beta_s = \mu_s/(\mu_s+\mu_p)$.
  \item \textbf{Khismatullin \& Nadim (2001)} \cite{khismatullin2001}:
        shape oscillations of a viscoelastic drop including the full
        Oldroyd/Jeffreys (relaxation + retardation) constitutive law;
        their characteristic equation already covers the Oldroyd-B case.
  \item \textbf{Zrni\'{c} \& Brenn (2024)} \cite{zirnic2024}:
        free-surface Oldroyd-B drop oscillations; their Eqs.\,(3.11)--(3.12)
        are the same characteristic equation and $q_*^2$ derived in \S\ref{sec:oldroydB}.
        The present notes expand the derivation and connect it to the
        Mola\v{c}ek \& Bush two-number framework.
  \item \textbf{Mukherjee \& Sarkar (2010)} \cite{mukherjee2009}:
        Oldroyd-B drop in a Newtonian bath; numerical results for
        oscillation frequency and damping as functions of $De_1$ and
        $\beta_s$, usable as benchmark for the free-surface limit.
\end{itemize}


\begin{thebibliography}{9}
\bibitem{reid1960}
  W.~H.~Reid,
  ``The Oscillations of a Viscous Liquid Drop,''
  \emph{Quart.\ Appl.\ Math.} \textbf{18}(1), 86--89 (1960).
\bibitem{molacek2012}
  J.~Mola\v{c}ek and J.~W.~M.~Bush,
  ``A quasi-static model of drop impact,''
  \emph{Phys.\ Fluids} \textbf{24}, 127103 (2012).
\bibitem{chandrasekhar1959}
  S.~Chandrasekhar,
  ``The oscillations of a viscous liquid globe,''
  \emph{Proc.\ London Math.\ Soc.} \textbf{9}(3), 141--149 (1959).
\bibitem{lamb1932}
  H.~Lamb,
  \emph{Hydrodynamics}, 6th~ed., Cambridge University Press (1932).
\bibitem{miller1968}
  C.~A.~Miller and L.~E.~Scriven,
  ``The oscillations of a fluid droplet immersed in another fluid,''
  \emph{J.\ Fluid Mech.} \textbf{32}, 417--435 (1968).
\bibitem{oldroyd1950}
  J.~G.~Oldroyd,
  ``On the formulation of rheological equations of state,''
  \emph{Proc.\ R.\ Soc.\ Lond.\ A} \textbf{200}, 523--541 (1950).
\bibitem{khismatullin2001}
  D.~B.~Khismatullin and A.~Nadim,
  ``Shape oscillations of a viscoelastic drop,''
  \emph{Phys.\ Rev.\ E} \textbf{63}, 061508 (2001).
\bibitem{mukherjee2009}
  S.~Mukherjee and K.~Sarkar,
  ``Effects of viscoelasticity on the retraction of a viscous drop,''
  \emph{J.\ Non-Newt.\ Fluid Mech.} \textbf{165}, 340--349 (2010).
\bibitem{bird1987}
  R.~B.~Bird, R.~C.~Armstrong, and O.~Hassager,
  \emph{Dynamics of Polymeric Liquids, Vol.\,1: Fluid Mechanics},
  2nd~ed., Wiley (1987).
\bibitem{zirnic2024}
  N.~Zrni\'{c} and G.~Brenn,
  ``Oscillations of a viscoelastic drop in the Oldroyd-B model,''
  \emph{Proc.\ R.\ Soc.\ A} \textbf{480}, 20230887 (2024).
\end{thebibliography}

\end{document}
